\section{Imbalance two: two returns of span three}\label{sec:tworeturns}
In this section, we prove that every orientation in the family~\ref{ov:family2} has a perfect fractional homomorphic tiling. As for internal roots in Section~\ref{sec:internalroots}, we fix an edge $x\to y$ of minimum weight and look for a pointwise inequality between the rescaled weights and a combination of the indicator functions of the neighbourhoods $X,Y,O,I$ of $x$ and $y$. The new difficulty is that our lower bounds on $X$ or $Y$ can now be smaller than $1/2$, and we compensate for this by bounding the number of vertices of small weight in $XI$ and in $YO$.

\begin{theorem}
\label{compact390:two-block}
Let $U$ and $V$ be closed
walks on $\{0,1,2,3\}$ which start and end at height zero and both visit height three. Then every digraph $F$ with $\delta^0(F)>|F|/2$ has a perfect fractional homomorphic $(+U+V)$-tiling.
\end{theorem}

\subsection{Neighbourhood bounds}
\label{compact390:maps-section}
Throughout this section, $H=+U+V$, where $U$ and $V$ are closed
walks on $\{0,1,2,3\}$ which start and end at height zero and
visit height three. Let $(a,b,c)$ and
$(d,e,f)$ be their crossing counts, as defined in Section~\ref{height390:three-section}, which are all positive, and let
\[
 B=a+d,\quad g=b+e,\quad T=c+f,\quad S=B+g+T+1,
 \qquad |H|=2S.
\]
We call $B$, $g$ and $T$ the \emph{bottom}, \emph{middle} and \emph{top} crossing counts of $H$; the bottom and top counts are its \emph{outer} crossing counts.
For each of $U$ and $V$, consider the steps before its first step from height two to height three, and count the steps from one down to zero and from two down to one among them. Let $R$ and $\iota$ be the totals of these two numbers over $U$ and $V$, so that
$0\le R\le B-2$ and $0\le\iota\le g-2$. As in Section~\ref{height390:three-section}, we write $PQ=P\cap Q$
for sets $P,Q$, and $P\ge c$ if $\omega(z)\ge c$ for every $z\in P$.

For the rest of the section, fix the following. Let $F$ be a digraph of order $N$ with minimum
semidegree $d_0>N/2$, and let $w$ be a deficient optimal fractional $H$-cover of $F$, so that
$\sum w<N/(2S)$ and $0\le w\le1$ by Lemma~\ref{general:LP}. Let $M=\max w$, let
$\kappa=\min_{xy\in E(F)}(w(x)+w(y))$, and let $\zeta=\max_{v,\sigma}\min_{N^\sigma(v)}w$ (this is the quantity called $\chi$ in Section~\ref{framework}).
Let $\omega=Sw$, so that every homomorphism from $H$ has $\omega$-cost at least $S$, and $\sum\omega<N/2$; all the weight bounds below are stated in terms
of $\omega$. Choose an edge $x\to y$ of minimum weight, and let $u=\omega(x)$, $v=\omega(y)$ and
$k=u+v$. Let
\[
 X=N^-(x),\quad Y=N^+(y),\quad O=N^+(x),\quad I=N^-(y),
\]
so that $z\in XY$ gives a directed triangle $x\to y\to z\to x$, and $z\in OI$
gives a path $x\to z\to y$. Finally, we define the following numbers and functions of $u$ and $v$, which are the lower bounds given by the next lemma:
\[
 D_X=2B+g+2,\quad D_Y=2T+g,\quad D=B+T+2g,
\]
\[
 A_R=2B+g-R-\iota,\quad D_R=D+R-\iota,\quad
 C_R=2T+g+2+R+\iota,\quad c_R=g+1+R-\iota,
\]
and
\[
\begin{split}
 F_X&=\frac{g+1+(B+T-g-1)u+(2B-1)v}{D_X},\qquad
 F_Y=\frac{g+1+(2T-3)u+(B+T-g-1)v}{D_Y},\\
 J_c&=\frac{g+1+(2T-3)u+(2B-1)v}{D},\qquad\quad\ \
 \beta_c=\frac{g-1+(T+1)u+(B+3)v}{D},\\
 F_R&=\frac{c_R+(B+T-g-1-2R+2\iota)u+(2B-3-2R)v}{A_R},\\
 J_R&=\frac{c_R+(2T-1+2\iota)u+(2B-3-2R)v}{D_R}.
\end{split}
\]
\begin{lemma}
\label{compact390:normalized-cover}\label{compact390:paired-floors}
With the notation above, the following hold.
\begin{props}
\item\label{compact390:normalized-cover:1} We have $\min w=0$, $2M+(S-2)\kappa\le1$, $k<1/2$, $2SM\le S-(S-2)k$ and $S\zeta<1-k$. Moreover, every $2$-cycle $p\leftrightarrow q$ satisfies
$\omega(p)+\omega(q)\ge1$.
\item\label{compact390:normalized-cover:2} We have the lower bounds
\[
 X\ge F_X,\quad Y\ge F_Y,\quad OI\ge J_c,
 \quad XY\ge1-\beta_c,\quad X\ge F_R,\quad OI\ge J_R.
\]
Moreover, each of the denominators above is positive, and $0<\beta_c<1/2$.
\end{props}
\end{lemma}
\begin{proof}
\emph{Part~\ref{compact390:normalized-cover:1}.} The sign sum of $H$ is two, $2S\ge14$, and its maximum drop is three, by Lemma~\ref{height400:last-exits}, as both returns have span three. Theorem~\ref{norm391:normalized-cover}
therefore gives the first inequalities, and Lemma~\ref{norm391:savings}
gives $k=S\kappa<1/2$, $S\zeta<1-k$ and the bound for $2$-cycles.

\emph{Part~\ref{compact390:normalized-cover:2}.} We use two partial homomorphisms. First, delete the initial position of $U$, between the separating edge before $U$
and the first step of $U$. The numbers of positions of the remaining path at its five
heights are
\[
 (a,\ B+b+1,\ c+d+g,\ T+e,\ f),
 \qquad\text{with both endpoints at height }1.
\]
Indeed, the contributions of $U$ and of the separating edge followed by $V$ are
$(a,a+b,b+c,c,0)$ and $(0,d+1,d+e,e+f,f)$.
The corresponding deletion in $V$ interchanges $(a,b,c)$ and
$(d,e,f)$, and so the sum of the two multiplicity vectors is $(B,\ D_X,\ D,\ D_Y,\ T)$.

Second, consider the first step of $U$ from height two to height three, and let $r$ and $i$ be the numbers of steps from one to zero and from two to one before it. This step is preceded by a step from one to two, and the part of $U$ before this step from one to two has
$(r+1,r+i+1,i,0)$ positions at heights $0,1,2,3$. Delete the middle
position of these two steps $1\to2\to3$. Read the rest of $U$ with its original heights, then
the separating edge and $V$ with heights increased by one, and then the initial part of $U$ with heights increased by two. Counting the positions gives
\[
 (a-r,\ B+b-r-i,\ c+d+g+r-i,
                     T+e+1+r+i,\ f+i),
 \qquad\text{with both endpoints at height }3.
\]
Again, every edge goes from one height to the next, and adding the multiplicity vector of the corresponding partial homomorphism
for $V$ gives $(B-R,\ A_R,\ D_R,\ C_R,\ T+\iota)$.

For the first pair of partial homomorphisms, send the heights one, two and three to
$(z,x,y)$, $(x,y,z)$ or $(x,z,y)$, for $z\in X$, $z\in Y$ or $z\in OI$, respectively.
For the second pair, we use the assignments for $X$ and $OI$ only. The required
edges exist by the definitions of the neighbourhoods. Send the heights zero
and four to inneighbours and outneighbours of minimum weight, which have $\omega$-weight less than $1-k$. Complete the deleted position by a vertex of weight at most $M$ forming a $2$-cycle with the common image of the endpoints, which exists as $2d_0>N$.
The two completed homomorphisms have total cost at least $2S$, and their two
completing vertices have total $\omega$-weight at most $2SM\le S-(S-2)k$.
If $E_*$ is the number of positions at heights zero and four in the two partial homomorphisms together, then the positions at heights one, two and three cost at least $S-E_*+(S-2+E_*)k$. Here $E_*=B+T$ for the first pair, and $E_*=B+T-R+\iota$ for the
second. Subtracting the contributions of $x$ and $y$ and dividing by the
multiplicity of $z$ gives $F_X,F_Y,J_c,F_R,J_R$.

For $z\in XY$, reduce the heights of the first pair modulo three and
send the residues zero, one and two to $x$, $y$ and $z$, respectively, using the triangle
$x\to y\to z\to x$. The total multiplicities of $x$ and $y$ are then
$B+2T+g$ and $2B+T+g+2$, and that of $z$ is $D$.
The endpoints have residue one, so the same completion bound
gives $\omega(z)\ge1-\beta_c$.

The bounds $R\le B-2$ and $\iota\le g-2$ imply that every denominator is positive.
Finally $0<\beta_c<1/2$: the numerator of $\beta_c$ is positive, and twice
that numerator minus $D$ is affine in $(u,v)$ and negative at the three
vertices of $u,v\ge0$, $u+v\le1/2$ \chk{9.H}.
\end{proof}

Choosing where each return descends gives lower bounds on two of the triple intersections, and slightly better lower bounds on $XO$ and $YI$ than those given by $2$-cycles, as follows.

\begin{lemma}
\label{compact390:actual-reserves}
Choosing the last step from height one to height zero in each return gives a homomorphism from $H$
with fibre sizes
\[
                     (2,B,B+g,T+g,T)
\]
into the path \(p_0\to\cdots\to p_4\) with shortcut
\(p_0\to p_2\). Consequently
\[
 XYO\ge Q_O:=
       \frac{S-(T+g+2)u-(B+T)v}{B+g}.
\]
Mapping one bipartition class of $H$ to a single vertex gives
\[
 XO\ge1-u+\gamma,\qquad YI\ge1-v+\gamma,
 \qquad \gamma=\frac{T}{2S-T}(1-k).
\]
In particular, $XYOI\ge Q_O$.
\end{lemma}
\begin{proof}
Apply Lemma~\ref{descent374:full-maps}, whose fibre sizes are
$(2,B,B+g,T+g,T)$, with images $(x,y,z,x,y)$ for $z\in XYO$,
so that the cover inequality gives $Q_O$.

Choose $L$ sinks of $H$ in one bipartition class of $H$, send every vertex
of the other class to $x$, these sinks to $y$,
and the other vertices of their class to a vertex $z\in XO$.
The edges at the chosen sinks are mapped to $x\to y$, and all other edges
to the $2$-cycle $x\leftrightarrow z$. Thus the
fibre sizes are $(S,L,S-L)$, and
\[
 \omega(z)\ge1-u+\frac{L}{S-L}(1-k).
\]
Similarly, take $K$ sources in one bipartition class. The reversed construction, with the other class at $y$, the chosen sources at $x$, and
the remaining positions at $z\in YI$, has fibre sizes
$(K,S,S-K)$, and gives the corresponding bound with $1-v$.
The $T$ positions at height three are sinks, and $H$ has as many sources as sinks,
so we may choose $L,K\ge T/2$. Since $H$ has two consecutive
forward edges, neither chosen set is a whole bipartition class.
Thus $S-L,S-K>0$, and, as $t/(S-t)$ is increasing in $t$, we obtain the bound with $\gamma$.

\end{proof}

The converse of $H$ again has the form $+U'+V'$, with the following crossing counts. We use it to exchange the roles of the bottom and top crossing counts.

\begin{lemma}
\label{compact390:literal-converse}
The converse $H^{\rm op}$, read backwards and suitably rotated, has the form $+U'+V'$ of $H$, with
\[
 (B',T',g',R',\iota')=(T+\iota,B-R,g+R-\iota,\iota,R).
\]
Applying this operation twice gives $H$ again.
Moreover, $F$ has a perfect fractional homomorphic $H$-tiling if and only if $F^{\rm op}$ has a perfect fractional homomorphic $H^{\rm op}$-tiling.
\end{lemma}
\begin{proof}
Write $U=P_U+Q_U$ and $V=P_V+Q_V$, where the displayed positive steps are the first steps from height two to height three. Thus $P_U$ and $P_V$ go
from zero to two without visiting three, and $Q_U$ and $Q_V$ go
from three to zero. If $\overleftarrow W$ denotes the word $W$ with the
order of its signs reversed, but the signs not negated, then the
converse of $H$, read backwards, has the cyclic word $+\bigl(\overleftarrow P_U+\overleftarrow Q_V\bigr) +\bigl(\overleftarrow P_V+\overleftarrow Q_U\bigr)$. In the first new return, the reversed part $\overleftarrow P_U$ has heights
$2-h$, the next positive step reaches three, and
the reversed part $\overleftarrow Q_V$ has heights $3-h$. Thus both new returns start and end at height zero and stay in $\{0,1,2,3\}$. If the numbers $R$ and $\iota$ for $U$ and $V$ separately are
$(R_U,I_U)$ and $(R_V,I_V)$, then counting the steps between consecutive heights gives the new
crossing counts
\[
 (f+I_U,e+R_U-I_V,d-R_V),\qquad
 (c+I_V,b+R_V-I_U,a-R_U).
\]
These are positive, as $R_U\le a-1$, $R_V\le d-1$, $I_U\le b-1$ and $I_V\le e-1$.
Before the first step from height two to height three in each new return, the numbers of steps from one to zero and from two to one are $(I_U,R_U)$ or
$(I_V,R_V)$. Adding gives the claimed values of $R'$ and $\iota'$.
The first steps from two to three of the new returns are the old separating edges, so applying the
operation twice gives $H$ again. The last assertion holds since reversing the edges of $H$ and $F$ preserves homomorphisms and their multiplicities.
\end{proof}


\subsection{Vertices of small weight and pointwise inequalities}
\label{compact390:compensation-section}
\subsubsection{Vertices of small weight}
The lower bounds of Lemma~\ref{compact390:normalized-cover} on $X$ and $Y$ may be less than $1/2$. To compensate, we bound the number of vertices of small weight in $XI$ and in $YO$. We use $X,Y,O,I$ as above, and write
\[
 k=u+v<1/2,\qquad s=1/2-k,\qquad h_c=1-k,
\]
\[
 Z_{XI}=\{z\in XI:\omega(z)<1/2\}\qquad\text{and}\qquad
 Z_{YO}=\{z\in YO:\omega(z)<1/2\}.
\]
These sets may be empty, and we only use bounds involving them when they are nonempty.

\begin{lemma}
\label{compact390:run-constraints}
The following hold.
\begin{props}
\item\label{run:disjoint} $Z_{XI}$ is disjoint from $O\cup Y$, and $Z_{YO}$ is disjoint from $I\cup X$.
\item\label{run:acyclic} $F[Z_{XI}]$ and $F[Z_{YO}]$ are acyclic.
\item\label{run:edges} If $a\in Z_{XI}$, $z\in O\cup Y$ and $z\to a$, then $\omega(a)+\omega(z)\ge3/2-k$. If $b\in Z_{YO}$, $z\in I\cup X$ and $b\to z$, then $\omega(b)+\omega(z)\ge3/2-k$.
\end{props}
\end{lemma}
\begin{proof}
Every directed run of $H$ has length at most four, by Lemma~\ref{compact390:balanced-template}. Recall that every $2$-cycle $p\leftrightarrow q$ has
$\omega(p)+\omega(q)\ge1$. As $u,v<1/2$, a vertex of $XO$ or $YI$ has weight more than $1/2$, which proves~\ref{run:disjoint}.

For~\ref{run:acyclic}, suppose that $a_0\to a_1\to a_2$ is a walk in $Z_{XI}$. Map all sources of $H$ to $a_0$,
all sinks to $y$, and the directed runs of lengths $1,2,3,4$ to
\[
 a_0y,\quad a_0xy,\quad a_0a_1xy,\quad a_0a_1a_2xy.
\]
respectively. All these edges exist, as $a_0,a_1,a_2\in XI$, so this is a homomorphism, and every image has weight
less than $1/2$, which contradicts the cover inequality. (If $a_2=a_0$, then $a_0\leftrightarrow a_1$ would contradict the bound for $2$-cycles.) Thus $F[Z_{XI}]$ has no walk with two edges, and in particular it is acyclic. Reversing all edges of $H$ and $F$ gives the same conclusion for $F[Z_{YO}]$.

For~\ref{run:edges}, we first construct a homomorphism from $H$ to an auxiliary orientation $H_*$.

Put $H_*=(++++---)^2$, with cyclic labels $0,\ldots,13$, and
map each separating edge, followed by the part of the next return up to its first visit
to three, by height, to the corresponding four-edge arm of $H_*$.
Map the suffix by height to its three-edge arm, read backwards
from that sink to the next source. The split images agree, so this defines $f:H\to H_*$ without
changing the order of the return steps.
Now consider the digraph with edges $a\to x\to y\to z\to a$ and $a\to y$, and the three maps $\phi_1,\phi_2,\phi_3$ from $H_*$ to it which send the labels $0,1,\ldots,6$, and also $7,8,\ldots,13$, to
\[
 a,x,y,z,a,z,y,\qquad y,z,a,x,y,a,z\qquad\text{and}\qquad z,a,x,y,z,y,a,
\]
respectively. Each $\phi_i$ is a homomorphism, and the three images of each
label are distinct; the same holds after composing with $f$.
Summing the three cover inequalities, each position of $H$ contributes the weights of three distinct vertices of $\{a,x,y,z\}$, which gives $3S\le2S\{\omega(a)+\omega(x)+\omega(y)+\omega(z)\}$. The same bound holds for $H^{\rm op}$: compose $f^{\rm op}$ with
the isomorphism $H_*^{\rm op}\to H_*$, $i\mapsto11-i\pmod{14}$.

Now take $a\in Z_{XI}$, $z\in O\cup Y$ with $z\to a$.
The vertices $a,x,y,z$ are distinct: $a\in\{x,y\}$ or $z=a$ would give
a loop, and $z\in\{x,y\}$ would give a $2$-cycle with $a$ whose two weights are below $1/2$.
If $z\in Y$, then $a\to x\to y\to z\to a$ and $a\to y$ are edges, and the inequality above gives
$\omega(a)+\omega(z)\ge3/2-k$.
If $z\in O$, we use only the edges $xy,ax,ay,xz,za$,
and map every sink of $H$ to $y$.
Each source of $H$ starts two directed runs, of lengths $r_1,r_2\in\{1,2,3,4\}$ say, and we map the source and these two runs as follows. If neither run has length two, we map the source to $x$, and a run of length one, three or four to $xy$, $xzay$ or $xzaxy$, respectively. If one of the runs has length two but neither has length three, we map the source to $a$, and a run of length one, two or four to $ay$, $axy$ or $axzay$. Finally, if the runs have lengths two and three, we map the source to $z$, and the runs to $zay$ and $zaxy$.
These walks have the required lengths, and they all end at $y$, so the images agree at the sinks.
Let $n_x,n_y,n_a,n_z$ be the numbers of positions mapped to $x,y,a,z$. We count them source by source, counting the source, the interiors of its two runs, and one sink (recall that $H$ has as many sinks as sources). Write $c_j$ for the number of runs of length $j$ among the two. In the first case the source contributes
$(1+c_4,1,c_3+c_4,c_3+c_4)$ to $(n_x,n_y,n_a,n_z)$, in the second case
$(c_2+c_4,1,1+c_4,c_4)$, and in the third case $(1,1,2,1)$.
Consequently their totals have
$n_a\ge n_z$, $n_x,n_y\ge1$, $n_z\le n_x+n_y$.

If $a_*=\omega(a)<1/2$ and
$\omega(z)<3/2-k-a_*$, this map's cost is at most
\[
 (n_a-n_z)a_*+\tfrac32n_z
                  +(n_x-n_z)u+(n_y-n_z)v
 \le\tfrac12\{n_a+n_z+\max(n_z,n_x,n_y)\}\le S.
\]
The first bound is strict if $n_z>0$; if $n_z=0$,
the last is strict since $n_x,n_y\ge1$. Thus its cost is
always below $S$, a contradiction.
Reversing all edges of $H$ and $F$, and applying the same argument to the edge $y\to x$ of $F^{\rm op}$,
gives the bound for edges from $Z_{YO}$ to $I\cup X$.
\end{proof}

Lemma~\ref{compact390:run-constraints} implies that there are few vertices of small weight, and that they can be matched to vertices of large weight, as follows. In it, let $Z$ and $\mathcal T$ be disjoint sets of vertices of $F$, such that $\omega<1/2$ on $Z$, $|\mathcal T|\ge d_0$ and $F[Z]$ is acyclic.
Suppose that either every edge from $\mathcal T$ to $Z$, or every edge from $Z$ to $\mathcal T$, has endpoints of total weight at least $3/2-k$.

\begin{lemma}
\label{compact390:cuts}\label{compact390:cuts:rest}
For $t\ge0$, whenever $Z_t=\{z\in Z:\omega(z)<1/2-t\}$ is nonempty,
\[
 |Z_t|+|\{w\in\mathcal T:\omega(w)<h_c+t\}|\le N-d_0.
\]
There is an injection $\pi:Z\to\mathcal T$ with $\omega(z)+\omega(\pi(z))\ge3/2-k$.
Consequently, for every $c\ge0$,
\[
 \sum_{z\in Z}\min(c,1/2-\omega(z))
 \le\sum_{w\in\mathcal T}\min(c,(\omega(w)-h_c)_+).
\]
\end{lemma}
\begin{proof}
In the first case, choose a vertex of $Z_t$ with no inneighbour in $Z_t$, which exists as $F[Z_t]$ is acyclic; in the second case, choose one with no outneighbour in $Z_t$.
Its inneighbourhood (respectively, outneighbourhood), which has at least $d_0$ vertices, avoids $Z_t$ and the vertices of $\mathcal T$ of weight less than $h_c+t$, as an edge between it and such a vertex would have endpoints of total weight
less than $1/2-t+h_c+t=3/2-k$. This gives the first inequality.

For the injection, put $e_z=1/2-\omega(z)>0$ and $c_w=(\omega(w)-h_c)_+$.
For $t>0$ with $\{e_z>t\}$ nonempty, the cardinality bound and $|\mathcal T|\ge d_0$ give $|\{c_w\ge t\}|\ge|\{e_z>t\}|+2d_0-N\ge|\{e_z>t\}|$. Letting $t$ increase to each demand $e$ and using finiteness
gives $|\{c_w\ge e\}|\ge|\{e_z\ge e\}|$.
Thus, matching the vertices of $Z$ in decreasing order of $e_z$ to the vertices of $\mathcal T$ in decreasing order of $c_w$ gives an injection with $c_{\pi(z)}\ge e_z$ for every $z$, that is, $\omega(z)+\omega(\pi(z))\ge3/2-k$. Summing $\min(c,e_z)\le\min(c,c_{\pi(z)})$ over $z\in Z$ gives the last inequality.
\end{proof}
By Lemma~\ref{compact390:run-constraints}, we may apply Lemma~\ref{compact390:cuts} to $Z=Z_{XI}$ and $\mathcal T=O\cup Y$, with edges into $Z$,
and to $Z=Z_{YO}$ and $\mathcal T=I\cup X$, with edges out of $Z$.
The two injections may overlap, and they need not follow edges.

\subsubsection{Pointwise inequalities with vertices of small weight}
We now give pointwise inequalities which also use the bounds of Lemma~\ref{compact390:cuts}. Throughout, $\omega\ge0$ is a weighting of a set of $N$ vertices, and $X,Y,O,I$ are sets of at least $d>N/2$ vertices each, with lower bounds
$l_X,l_Y,v,u$, respectively, where $u,v\ge0$, $k=u+v<1/2$,
$s=1/2-k$ and $h_c=1-k$. Suppose that
\[
 XY\ge1-\beta,\quad OI\ge J,\quad
 XO\ge1-u+\Gamma,\quad YI\ge1-v+\Gamma,\quad \Gamma\ge0.
\]
Let $a_s=(1/2-l_X)_+$ and $b_s=(1/2-l_Y)_+$, and let
\[
 Z_{XI}=\{z\in XI:\omega(z)<1/2\},\quad E_{XI}=Z_{XI}\cup\{O\cup Y:\omega<h_c\},
\]
\[
 Z_{YO}=\{z\in YO:\omega(z)<1/2\},\quad E_{YO}=Z_{YO}\cup\{I\cup X:\omega<h_c\}.
\]
In the next lemma, the function $P$ is a combination of the indicator functions of $X,Y,O,I$, from which we subtract $a\ind_{E_{XI}}$ and $b\ind_{E_{YO}}$ to account for vertices of small weight in $XI$ and $YO$. By Lemma~\ref{compact390:cuts}, the sets $E_{XI}$ and $E_{YO}$ have at most $N-d$ vertices when $Z_{XI}$ and $Z_{YO}$ are nonempty. A further correction, with coefficient $\mu$, uses the last inequality of Lemma~\ref{compact390:cuts}.
In Lemma~\ref{short392:incidence-table}, we fix parameters $D,e_A,e_B\ge0$, $0\le a,b\le s$, $0\le\mu\le1$ and $c\ge0$, and define
\[
\begin{gathered}
 c_I=u+D+b-e_B,\quad c_O=v+D+a-e_A,\quad
 E=e_A+e_B,\quad \Lambda=a+b,\quad \eta=\mu c,\\
 C(q)=\mu\min(c,q),\quad
 q_0=\max(0,k-\beta,l_X-h_c,l_Y-h_c),\quad
 \mathcal B=D+\Lambda+\eta.
\end{gathered}
\]

\begin{lemma}
\label{short392:incidence-table}
In the setting above, suppose that $|E_{XI}|\le N-d$ if $Z_{XI}$ is nonempty and $|E_{YO}|\le N-d$
if $Z_{YO}$ is nonempty, that $a=0$ if $Z_{XI}$ is empty and $b=0$ if $Z_{YO}$ is empty, and, if $\mu>0$, that
\[
 \sum\left[(1/2-\omega)_+\ind_{YO}
       -\min(c,(\omega-h_c)_+)\ind_{I\cup X}\right]\le0.
\]
Suppose also that
\[
\begin{gathered}
 c_I\ge a_s,\quad c_O\ge b_s,\quad
 a\ge a_s\ (Z_{XI}\ne\varnothing),\quad
 b+\mu b_s\ge b_s\ (Z_{YO}\ne\varnothing),\\
 \beta-k\le D+\Lambda-E\le2s,\quad k+D-E\le J,\quad
 q_0-C(q_0)\ge E-D,\quad \mathcal B\le s+\Gamma.
\end{gathered}
\]
and that the cyclic triple intersections and $XYOI$ satisfy
\[
 XYO\ge1-u+a+e_B+\eta,\quad
 XYI\ge1-v+b+e_A+\eta,\quad XYOI\ge1+\mathcal B.
\]
Then $\sum\omega>N/2$, whether the sets $Z_{XI}$ and $Z_{YO}$ are empty or not.
\end{lemma}
\begin{proof}
Set $c_X=1/2-u+e_B$, $c_Y=1/2-v+e_A$ and
\[
 P=\sum_Rc_R\ind_R-D-a\ind_{E_{XI}}-b\ind_{E_{YO}}
 -\mu(1/2-\omega)_+\ind_{YO}+C((\omega-h_c)_+)\ind_{I\cup X}.
\]
We check that $P\le\omega$ at every vertex, according to which of $X,Y,O,I$ contain it. All the coefficients $c_R$ are nonnegative. On $XO$ and $XOI$, dropping the negative terms, using $C\le\eta$,
and adding $c_I$ if necessary, we see that $P$ is at most $1-u+\mathcal B-s-e_A\le1-u+\Gamma$, which is a lower bound on $XO$.
Similarly, $P\le1-v+\mathcal B-s-e_B\le1-v+\Gamma$
on $YI$ and $YOI$.
On $XYO$, $XYI$ and $XYOI$,
the same sums of coefficients give exactly the three assumed lower bounds.

Consider now the remaining patterns at a vertex with $\omega<h_c$. The values of $P$ on $X,Y,O,I$ are
$1/2-c_I$, $1/2-c_O$, $v-e_A$ and $u-e_B$, which are at most $l_X,l_Y,v,u$, respectively.
The values on $OI$ and $XY$ are $k+D-E\le J$ and
$h_c-D-\Lambda+E\le1-\beta$.
A vertex of $XI$ or $YO$ of weight less than $1/2$ lies in no other of the four sets,
by the lower bounds on $XO$ and $YI$. For such a vertex, let $e=1/2-\omega$, so that $e\le a_s$ or $e\le b_s$, respectively;
then $P-\omega$ is $e-a$ or $(1-\mu)e-b$, respectively, which is
nonpositive by our assumptions, as the corresponding set is nonempty.
For $1/2\le\omega<h_c$, the values on $XI$ and $YO$ are both $1/2$.

Now let $\omega\ge h_c$, and put $q=\omega-h_c$. Without the correction term $C$,
the values on $XI$, $YO$ and $OI$ are $1/2+b$, $1/2+a$ and $k+D+\Lambda-E$, which are
all at most $h_c$. As the coefficients are nonnegative, the values on $X$, $Y$, $O$ and $I$ are at most these, and $C(q)\le q$ gives $P\le\omega$.
On $XY$, we have $q\ge q_0$ and
$P-\omega=E-D+C(q)-q\le0$, since $q-C(q)$ is nondecreasing
and $q_0-C(q_0)\ge E-D$. Finally, outside the four sets $P=-D$. This covers all sixteen patterns.

The coefficients $c_R$ sum to $1+2D+\Lambda$. Summing $P\le\omega$ over all vertices, and
using the cardinality bounds for the nonempty sets among $Z_{XI},Z_{YO}$, and the inequality assumed when $\mu>0$, gives
\[
 \sum\omega\ge(1+2D+\Lambda)d-DN-\Lambda(N-d)
          \ge N/2+(1+2D+2\Lambda)(d-N/2)>N/2.
\]
If $Z_{XI}$ or $Z_{YO}$ is empty, the corresponding coefficient $a$ or $b$ is zero, so no cardinality bound is needed for it.
\end{proof}

We will use two choices of coefficients in Lemma~\ref{short392:incidence-table}, as follows. The first is for the case where the lower bounds on $X$ and $Y$ are at least $k$, and the second for the case where the lower bound on $X$ is at least $1/2$ but that on $Y$ is less than $k$.

\begin{lemma}
\label{short392:incidence-table-two}\label{compact390:regular-certificate}\label{compact390:weak-certificate}
In the setting above, suppose that $0\le\beta<1/2$ and that the conclusions of
Lemma~\ref{compact390:cuts} hold for $Z_{XI}$ and $Z_{YO}$, when they are nonempty. Let
\[
 \tau=\max(0,a_s+b_s-s,2a_s-u-s,2b_s-v-s).
\]
Suppose also that $XYI\ge1-v+b_s$.
Then either of the following sets of conditions gives $\sum\omega>N/2$.
\begin{props}
\item\label{short392:incidence-table:R} $l_X,l_Y\ge k$, $J\ge\max(k,\beta,v+a_s)$, $\Gamma=\tau$,
\[
 XYO\ge1-u+s+\tau,\qquad XYOI\ge1+s+\tau.
\]
\item\label{short392:incidence-table:W} $l_X\ge1/2$, $u\le l_Y<k$, $\beta<k$,
$J\ge\max(u,v,\beta)$, $\Gamma\ge b_s$, and
\[
                  XYO,XYOI\ge1+s+b_s.
\]
and, moreover,
\[
 q=\max(k-\beta,l_X-h_c)\ge\chi,\qquad\chi=e+t\min(1,e/s),
 \qquad e=(k-J)_+,\quad t=k-l_Y.
\]
\end{props}
Note that, in~\ref{short392:incidence-table:W}, the set $Z_{XI}$ is empty, as $l_X\ge1/2$, and the bound on $XYI$ follows from the bound on $YI$.
\end{lemma}
\begin{proof}
To apply Lemma~\ref{short392:incidence-table}, put $t=(b_s-s)_+$ and take
\[
 a=a_s\ind_{Z_{XI}\ne\varnothing},\quad
 b=\min(b_s,s)\ind_{Z_{YO}\ne\varnothing},\quad
 \mu=(t/b_s)\ind_{Z_{YO}\ne\varnothing},\quad c=b_s,
\]
with $\mu=0$ if $b_s=0$. Define
\[
 D=\max(0,a_s-u-b,b_s-v-a,\beta-k-a-b),\quad
 e_A=0,\quad e_B=E=(k+D-J)_+,\quad \eta=\mu b_s.
\]
We use the displayed formula for $\tau$ in both cases. In \ref{short392:incidence-table:R} we have $t=\eta=0$ and $\tau\le s$, whereas in \ref{short392:incidence-table:W} $a_s=0$,
$b_s=s+t\le s+v$, $D=(b_s-v)_+\le s$ and $\tau=D+t\le b_s$.
Thus $a_s\le s$, $b+\eta\le b_s$, $\eta\le t\le\tau$ and $a_s\eta=0$.
Consequently $\mathcal B=D+a+b+\eta \le\max(a_s+b_s,2a_s-u,2b_s-v,s+\eta,\eta-u)\le s+\tau$. This bounds the four defining entries of $D+a+b$, using $a_s\eta=0$.
Both sets of hypotheses give $J\ge\max(u,v+a_s,\beta)$; hence
\[
 c_I=\min(u+D+b,J-v+b)\ge a_s,\quad c_O=v+D+a\ge b_s,
\]
and $D+a+b-E=\min(D+a+b,J-k+a+b)\ge\beta-k$.
Its upper bound $2s$ follows from $\tau\le s$ in \ref{short392:incidence-table:R}, or $D,b\le s$
in \ref{short392:incidence-table:W}; also $k+D-E\le J$.
For nonempty sets, $a=a_s$ and $b+\mu b_s=b_s$, which gives the
conditions on $a$ and $b$. If $\mu>0$, then $Z_{YO}$ is nonempty and its
demands are at most $b_s$, so the inequality with $c=b_s$ applies.
An empty set has zero coefficient in the cardinality bound; empty $Z_{YO}$ also has $\mu=0$.

Since $E-D\le e=(k-J)_+$, it remains to show $q_0-C(q_0)\ge e$.
For $\mu=0$ this follows from $e=0$ in \ref{short392:incidence-table:R}, or $q_0=q\ge\chi\ge e$ in \ref{short392:incidence-table:W}.
For $\mu>0$ only \ref{short392:incidence-table:W} occurs. The function $q'-\mu\min(b_s,q')$
equals $(s/b_s)q'$ up to $b_s$ and $q'-t$ thereafter; its inverse at $e$ is
\[
 \begin{cases}e+te/s,&e\le s,\\e+t,&e\ge s,\end{cases}
                  =e+t\min(1,e/s)=\chi.
\]
Thus the exact threshold holds, including $e=0$.

The pair, fourfold and $XYI$ bounds follow from $\Gamma\ge\tau$,
$\mathcal B\le s+\tau$ and $b+\eta\le b_s$.
For $XYO$, use the common estimate
$a+E+\eta\le\mathcal B+e\le s+\tau+e$.
In \ref{short392:incidence-table:R} we have $e=0$, and in \ref{short392:incidence-table:W} $e\le u$ and $\tau\le b_s$, so this gives both stated bounds. Lemma~\ref{short392:incidence-table}
applies, including empty sets; its inequality is used only when $\mu>0$.
\end{proof}

\subsubsection{Comparing the lower bounds}
To apply these inequalities, we need some comparisons between the lower bounds in Lemma~\ref{compact390:normalized-cover}\ref{compact390:normalized-cover:2}. We write $n=\min(B,T)$ and $h=|B-T|$ for the smaller outer crossing count and the difference of the outer crossing counts.

\begin{lemma}
\label{bulk374:actual-inputs}
\label{inherit376:actual-cut}
Let $n=\min(B,T)\ge2$, $h=|B-T|$, $g\ge2$, and
$D=2n+h+2g$. Let $\varepsilon=1$ and $t=v$ if $B\ge T$,
and $\varepsilon=0$ and $t=u$ otherwise, so that $h\ge1$ in the latter case.
The lower bounds of Lemma~\ref{compact390:normalized-cover}\ref{compact390:normalized-cover:2} on $X$ and $Y$, on $OI$ and on $XY$ can be written as
\[
\begin{split}
 L&=\frac{g+1+(2n+h-g-1)k+(h+g-2+2\varepsilon)t}
              {E},\qquad E=2n+2h+g+2\varepsilon,\\
 L^*&=\frac{g+1+(2n-1-2\varepsilon)k+(h-g+2\varepsilon)t}
              {2n+g+2-2\varepsilon},\\
 J_c&=\frac{g+1+(2n-1-2\varepsilon)k+(2h-2+4\varepsilon)t}{D},\\
 \beta&=\frac{g-1+(n+3-2\varepsilon)k+(h-2+4\varepsilon)t}{D}.
\end{split}
\]
where $L$ and $L^*$ are the lower bounds on $X$ and $Y$ if $\varepsilon=1$, and on
$Y$ and $X$ if $\varepsilon=0$, and $J_c$ and $1-\beta$ are the lower bounds on $OI$ and $XY$.
They satisfy
\[
 L^*>k,\quad L>t,\quad L>k/2,\quad J_c>t,\quad J_c>\beta,
 \quad0<\beta<1/2,
\]
and $\beta\ge k$ implies $L>k$.
The bound $J=\max(J_c,u,v)$ further satisfies
\[
 L+J\ge k+t,\qquad
 j=(k-J)_+,\ a_0=(k-L)_+:\quad L\ge t+j,\quad j+a_0\le k-t.
\]
\end{lemma}
\begin{proof}
Substituting $B$ and $T$ into the bounds of
Lemma~\ref{compact390:normalized-cover}\ref{compact390:normalized-cover:2}
gives the four expressions. Every comparison in the first display is affine
on the triangle $0\le t\le k\le1/2$, so it suffices to check the
vertices \chk{9.I}. There $L^*-k$, $L-t$ and $J_c-t$ are nonnegative and
positive at the origin, hence strict for $k<1/2$, while $L-k/2$,
$J_c-\beta$, $\beta$ and $1/2-\beta$ are positive at all three vertices;
also $\beta\le\bar\beta(k)=(g-1+(n+h+3)k)/D$.
The polygon $\beta\ge k$ is spanned by the origin and its intercepts with
$t=0$ and $t=k$, and $L-k$ is positive at all three points \chk{9.I}.

To prove the bound involving $J$, put
\[
 A=(g+1)(D+E),\quad b=(h+2g+1+2\varepsilon)(D+E)-DE,
 \quad Z=b-2A,\quad c=2Z+D^2,
\]
and $f=DE(L+J_c-k-t)=A-bk+ct$. If $Z\le0$, this affine function is
nonnegative at all three vertices. If $Z>0$, then $c>0$. Put
\[
 a_t=4n+3h+g+2+2\varepsilon>0,\qquad
 p=E(L-2t)=g+1+(2n+h-g-1)k-a_t t.
\]
The positive combination $a_tf+cp$ does not depend on $t$, and it is
positive at $k=0$ and at $k=1/2$ \chk{9.I}. Thus $f<0$ forces $p>0$.
Since $J\ge k-t$, this gives $L+J-k-t\ge L-2t>0$, and hence
$L+J\ge k+t$. If both $j,a_0$ are positive this yields
$j+a_0\le k-t$, while if only one is positive we use $L>t$ or $J>t$;
the bound $L\ge t+j$ follows in the same way.
\end{proof}

\subsection{A sufficient condition}
Suppose that the top crossing count exceeds the bottom one by more than the middle count. Then the conditions of Lemma~\ref{short392:incidence-table-two} follow from a single inequality in $u$ and $v$, together with two inequalities in the crossing counts, as follows. This inequality is verified in Section~\ref{sec:remaining}.
\begin{proposition}
\label{compact390:source-criterion}
Let $H=+U+V$ have $B=n\ge2$, $T=n+h$ and $n\le g<h$.
Let $R$ and $\iota$ be as defined at the start of this section,
and let $F_X,F_Y,J_c,\beta_c,F_R$ be as in
Lemma~\ref{compact390:normalized-cover}\ref{compact390:normalized-cover:2}, where now
\[
 D_X=2n+g+2,\quad E=2n+2h+g,\quad D=2n+h+2g,
\]
Suppose that
\[
 N_1=2(g+1)(D+D_X)-DD_X>0,\qquad
 N_X=2g^2+(h-1)(g-2n)-2\ge0.
\]
Let $s=1/2-k$, $h_c=1-k$ and $\overline J=\max(J_c,u,v)$, and
suppose that, for all $u,v\ge0$ with $k=u+v<1/2$,
$F_Y<k$ and $\overline J<k$, we have
\[
 \max\{k-\beta_c,F_X-h_c,F_R-h_c\}
 \ge e+t\min(1,e/s),\qquad e=k-\overline J,\quad t=k-F_Y,
\]
Then every digraph $F$ with $\delta^0(F)>|F|/2$ has a perfect fractional homomorphic $H$-tiling.
\end{proposition}
\begin{proof}
Suppose that a deficient optimal cover exists, and use the notation and bounds of Lemma~\ref{compact390:normalized-cover}. By Lemma~\ref{compact390:run-constraints}, the conclusions of Lemma~\ref{compact390:cuts} hold for $Z_{XI}$ and $Z_{YO}$.
Applying Lemma~\ref{bulk374:actual-inputs} with
$\varepsilon=0$ and $t=u$ identifies its $L^*,L,J_c,\beta$ with
$F_X,F_Y,J_c,\beta_c$, so that
$F_X>k$, $F_Y>u$, $J_c>\max(u,\beta_c)$ and
$0<\beta_c<1/2$.
In addition $F_X>F_Y$ \chk{9.J}.
Put $J=\overline J$, $a_s=(1/2-F_X)_+$ and $b_s=(1/2-F_Y)_+$.

Lemma~\ref{compact390:actual-reserves} gives
\[
 \gamma=\frac{T}{K}(1-k),\quad T=n+h,\quad K=2n+T+2g+2,
 \qquad XO\ge1-u+\gamma,\quad YI\ge1-v+\gamma,
\]
\[
 XYO\ge Q_O=\frac{2n+h+g+1-(n+h+g+2)u-(2n+h)v}{n+g}.
\]
On the triangle $u,v\ge0$, $u+v\le1/2$ we have
\[
 \gamma>b_s,\qquad Q_O\ge1+s+\frac{n+h-g}{2(n+g)},\qquad Q_O+F_Y>2-k;
\]
each is an affine comparison, checked at the vertices \chk{9.J}. (At the
$u$ half-axis $F_Y>1/2$, so there $b_s=0$.) Thus $Q_O>1+s+b_s$, which holds on $XYOI$ as well, and
$XYI\ge1-v+\gamma$ follows from $YI$.

Writing $a=h-1$, $L=2g+h+1$ and $A_0=g+h-2$, the definitions give
\[
 D(J_c-k)=g+1-Lk+2au,\qquad
 E(F_Y-k)=g+1-Lk+A_0u.
\]
Since $A_0<2a$, we have $J_c\ge k$ wherever $F_Y\ge k$.
We consider the two regions $F_Y\ge k$ and $F_Y<k$ of the triangle $u,v\ge0$, $u+v\le1/2$ separately, and we use condition~\ref{short392:incidence-table:R} of Lemma~\ref{short392:incidence-table-two} in the first and condition~\ref{short392:incidence-table:W} in the second. The two regions share the boundary vertices
\[
 P=(0,(g+1)/L),\qquad
 W=(a/(2A_0),1/2-a/(2A_0)).
\]
The first region is the quadrilateral with vertices $(0,0),(1/2,0),P,W$.
The affine comparison $J_c+F_X-1/2-v$ has first three values
$N_1/(2DD_X)$, $1/2+a/(2D)+a/(2D_X)$ and $N_X/(2LD_X)$ \chk{9.J}.
They are nonnegative by $N_1>0$ and $N_X\ge0$. At $W$,
$F_Y=1/2$, so $F_X>1/2$ and $J_c\ge1/2$. Hence $J\ge v+a_s$.
Since $a_s\le b_s\le s$ we also have $\tau\le b_s$. With
$l_X=F_X,l_Y=F_Y,\beta=\beta_c$, the bounds above then give all other
assumptions of \ref{short392:incidence-table:R} of Lemma~\ref{short392:incidence-table-two}.

The closure of the region $F_Y<k$ is the triangle $P,(0,1/2),W$. At $P$,
$F_X-1/2=N_X/(2LD_X)\ge0$; at the $v$ half-axis it is positive, and at $W$ we already have $F_X>1/2$, so
$F_X\ge1/2$ throughout.
We also have $\beta_c<k$, since its first vertex satisfies $J_c=k>\beta_c$
and both half vertices have $\beta_c<1/2$.
Set $l_X=\max(F_X,F_R)$, $l_Y=F_Y$, $J=\overline J$.
Then $l_Y>u$, $J\ge\max(u,v,\beta_c)$, and the threshold is $q=\max(k-\beta_c,l_X-h_c) =\max(k-\beta_c,F_X-h_c,F_R-h_c)$. If $J\ge k$, then $e=0$; otherwise the assumed inequality gives
$q\ge\chi$. All other conditions of \ref{short392:incidence-table:W} of Lemma~\ref{short392:incidence-table-two} were proved above.
In both regions, Lemma~\ref{short392:incidence-table-two} gives $\sum\omega>N/2$, which contradicts $\sum w<N/(2S)$. Thus there is no deficient cover, and
Lemma~\ref{general:LP} proves the proposition.
\end{proof}


\subsection{Large outer crossing counts}
\subsubsection{Bounds from directed runs}
\label{compact390:bulk-section}
When the outer crossing counts are large compared with the middle one, the lower bounds on the triple intersections given by Lemma~\ref{geom368:unrestricted-folds} are large, and a simpler pointwise inequality suffices.
We continue to use the cover $w$ of Section~\ref{compact390:maps-section}.


\begin{corollary}
\label{geom368:canonical-fold-budget}
For $H=+U+V$ with crossing counts $(a,b,c),(d,e,f)$,
each of the four triple intersections has lower bound
\[
                         \frac{S(1-k)}{2(b+e)+2},
\]
and $XYOI$ has lower bound this quantity plus $k/2$.
\end{corollary}
\begin{proof}
Let $A=\sum_h\lfloor j_h/2\rfloor$, where the sum is over the lengths $j_h$ of the directed runs of $H$. We assign the contribution of each run to steps of the returns between heights one and two, or to separating edges, so that distinct runs are assigned distinct steps. A run which does not
contain a separating edge has length at most three. If its
length is two or three, it contains a step between heights one and two, and it contributes one to $A$; we assign it to this step.
A run which contains a separating edge consists of positive steps. If its length is
two, it contributes one, which we assign to the separating edge. If its
length is three, its two steps in the following return contain the step from one to two, and its contribution one is assigned
to it. If its length is four, its three steps in the following return contain the step from one to two, and its contribution two is assigned
to this step and to the separating edge. A run of length one contributes nothing.
There are $2b+2e$ steps between heights one and two and two separating edges,
so $A\le2(b+e)+2$, and Lemma~\ref{geom368:unrestricted-folds} gives the bounds.
\end{proof}

Here we use the bounds of Lemma~\ref{compact390:cuts} at all thresholds $t$ at once, by integrating over $t$.

\begin{lemma}
\label{simpleunion:singleton-credit}
\label{top372:multilayer}
Let $\omega\ge0$ be a weighting of a set of $N$ vertices, and let $X,Y,O,I$ be sets of at least $d>N/2$ vertices. Let $u,v\ge0$, $k=u+v<1/2$,
$s=1/2-k$ and $h_c=1-k$. Suppose that $X,Y,O,I$ have lower bounds
$l_X,l_Y,v,u$, respectively, that $XO\ge1-u$ and $YI\ge1-v$,
and that $XY\ge1-\beta$ and $OI\ge J$, where
\[
 0\le\beta<1/2,\quad J\ge\beta,\quad
 j=(k-J)_+,\quad a_0=(k-l_Y)_+,\quad
 l_X\ge k,\quad l_Y\ge u+j,\quad j+a_0\le v.
\]
Suppose that, for every $t\ge0$, we have
\[
 \left|\{XI:\omega<1/2-t\}\cup\{O\cup Y:\omega<h_c+t\}\right|
 \le N-d,
\]
whenever $\{XI:\omega<1/2-t\}$ is nonempty, and the corresponding bound for $YO$ and $I\cup X$ whenever $\{YO:\omega<1/2-t\}$ is nonempty. If $\beta\le k$, suppose that
\[
 XYO\ge1+s-u,\quad XYI\ge1+s-v+j+a_0,\quad XOI\ge1+s-u-j+a_0,\quad YOI\ge1+s-v+a_0,
\]
and $XYOI\ge1+2s+a_0$.
If $\beta>k$, suppose instead that $l_Y\ge k$, that every triple intersection has lower bound
$1+s$, and that $XYOI$ has lower bound
\[
               1+\max\{2s,(\beta-2\min(u,v))_+\}.
\]
Then $\sum\omega>N/2$.
\end{lemma}
\begin{proof}
Let $Z_{XI}=\{z\in XI:\omega(z)<1/2\}$ and $Z_{YO}=\{z\in YO:\omega(z)<1/2\}$, and let
$E_{XI}=Z_{XI}\cup\{z\in O\cup Y:\omega(z)<h_c\}$ and
$E_{YO}=Z_{YO}\cup\{z\in I\cup X:\omega(z)<h_c\}$.
Let $a=(k-\min_{YO}\omega)_+$, with $a=0$ if $YO$ is empty;
thus $0\le a\le a_0$. For $0\le t<a$, the set
$\{z\in YO:\omega(z)<k-t\}$ is nonempty, so the bound for $YO$ at the threshold $s+t$ applies;
this set is disjoint from $I\cup X$, as vertices of $YI$ and $XO$ have weight at least $1/2$.
Integrating these bounds over $0\le t<a$ gives
\[
 \sum_z\left[(k-\omega(z))_+\ind_{YO}
       +\{a-(\omega(z)-h_c-s)_+\}_+\ind_{I\cup X}\right]
                         \le a(N-d).
\]
Indeed, a vertex
of $YO$ contributes the length $(k-\omega)_+\le a$ of the set of $t$ for which it is counted,
and a vertex of $I\cup X$ contributes $\{a-(\omega-h_c-s)_+\}_+$.

Let $(c_X,c_Y,c_O,c_I)=(1/2-u,1/2-v+j,v+s-j,u+s+a)$, which are nonnegative, and let $P$ be $\sum_Rc_R\ind_R$ minus $s\ind_{E_{XI}}+s\ind_{E_{YO}}$ and minus the integrand above. We check that $P\le\omega$ on each membership pattern.
For a vertex of weight less than $h_c$, the values of $P$ on $X,Y,O,I$ are at most $v,u+j,v-j,u$, which are at most
the corresponding lower bounds. For a vertex of weight at least $h_c$, the values on $X,Y,O$ are
at most $1/2$. With $q=\omega-h_c\ge0$, the $I$ value is
$u+s+r(q)$, where $r(q)=\min(a,(q-s)_+)$. If $q\le s$
this is at most $1/2$; otherwise it is at most $u+q\le\omega$.

A vertex of $XI$ or $YO$ of weight less than $1/2$ lies in no other of the four sets, by the bounds on $XO$ and $YI$.
On $XI$, the values of $P$ for weights in $[0,1/2)$, in $[1/2,h_c)$ and in $[h_c,\infty)$ are
$k$, $1/2$ and $h_c+r(q)$; the last is at most $\omega$, as $r(q)\le q$.
The corresponding values on $YO$ are
$k-(k-\omega)_+$, $1/2$ and $h_c$, which are at most $\omega$.
For $OI$ the lower value is $k-j\le J$, and the upper value
is $h_c-j+r(q)\le\omega$.
For $XY$, the lower value is at most $k+j\le1-\beta$.
Above $h_c$ its value is at most $h_c+j$: if $j>0$ this is
$1-J\le1-\beta$, and if $j=0$ it is at most the weight itself.
Finally, the values on $XO$ and $YI$ are at most
$1-2u-j$ and $1-2v+j+a$, which are at most $1-u$ and $1-v$ as $j+a\le v$.

Dropping nonpositive terms, the values on the four triple intersections and on
$XYOI$ are at most the five lower bounds stated above, with $a$ in place of $a_0$. When $\beta>k$, we have $j=a_0=a=0$, so these values are at most $1+s$ on the triple intersections and $1+2s$ on $XYOI$. Outside the four sets, $P$ is at most zero. Thus $P\le\omega$ in all sixteen cases.

We cannot use the cardinality bound for an empty set, so we modify $P$ as follows. If $Z_{XI}$ is empty, we remove the term $-s\ind_{E_{XI}}$, increase $c_X$ by $s$,
and subtract the constant $s$. The change
$s(\ind_X-1+\ind_{E_{XI}})$ is nonpositive for weights at least $h_c$.
For smaller weights, $XO$ cannot occur by its lower bound, and the only
pattern where the change is positive is $XY$, where the new value is
at most $1/2+j\le1-\beta$. If $Z_{YO}$ is empty, we remove the term $-s\ind_{E_{YO}}$, increase $c_Y$ by $s$,
and subtract $s$; in this case $a=0$. If both sets are empty,
the new value on $XY$ is $h_c+j$, which is still at most the lower bound when $\beta\le k$
or $j>0$. If $z\in\{0,1,2\}$ of the two sets are empty, then the sum of the coefficients, the
constant, and the total coefficient of the terms bounded by Lemma~\ref{compact390:cuts} are, respectively,
$1+2s+a+zs$, $zs$ and $(2-z)s+a$. Summing $P\le\omega$ gives
\[
 \sum\omega\ge N/2+(1+4s+2a)(d-N/2)>N/2.
\]

It remains to treat the case where $\beta>k$ and both sets are empty, which needs a different choice.
Here $j=a=0$ and the lower bounds on $X$ and $Y$ are at least $k$.
Put $m=\min(u,v)$ and
\[
 D=\max(0,2s-\beta,\beta-2m),\quad
 c_X=c_Y=1/2+(D-\beta)/2,\quad c_O=c_I=(D+\beta)/2.
\]
We use $P=\sum c_R\ind_R-D$, without any cardinality bound. Its values on $X$ and $Y$
are $1/2-(D+\beta)/2\le k$, and its values on $O$ and $I$ are $(\beta-D)/2\le m$.
Its values on $XI$, $YO$, $XO$ and $YI$ are $1/2$, which is at most $\omega$, as $Z_{XI}$ and $Z_{YO}$ are empty and by the lower bounds on $XO$ and $YI$. Its values on $XY$ and $OI$ are $1-\beta$ and $\beta$.
Its values on the cyclic triple intersections are $1+(D-\beta)/2\le1+s$, and on the other two triple intersections they
are at most one, since $D+\beta\le2\max(s,\beta)\le1$.
Its value $1+D$ on $XYOI$ is at most the assumed lower bound. The coefficients
are nonnegative and sum to $1+2D$, so
\[
 \sum\omega\ge(1+2D)d-DN
             =N/2+(1+2D)(d-N/2)>N/2.
\]
This completes the proof.
\end{proof}

\begin{theorem}
\label{compact390:bulk-finite}\label{compact390:outer-finite}
Let $n=\min(B,T)$ and $h=|B-T|$, so that $S=2n+h+g+1$. Then every digraph $F$ with $\delta^0(F)>|F|/2$ has a perfect fractional homomorphic $H$-tiling
if either of the following holds:
\begin{props}
\item\label{compact390:bulk-finite:1} \[
 S\ge4(g+1),\qquad
 R_b=S(n+g-2)-(g+1)(2n+5g-7)\ge0,
\]
\item\label{compact390:bulk-finite:2} $B+T\ge4g+1$.
\end{props}
\end{theorem}
\begin{proof}
\emph{Part~\ref{compact390:bulk-finite:1}.} Suppose that a deficient optimal cover exists. We use the bounds of
Lemma~\ref{compact390:normalized-cover}\ref{compact390:normalized-cover:2}, replacing $J_c$ by
$J=\max(J_c,u,v)$ as in Lemma~\ref{inherit376:actual-cut}.
Let $A_f=\sum\lfloor j/2\rfloor$, over the lengths $j$ of the directed runs of $H$, let $E_f$ be the number
of directed runs of even length, and let $\theta=E_f/(2A_f)\in[0,1/2]$.
Lemma~\ref{geom368:unrestricted-folds} and
Corollary~\ref{geom368:canonical-fold-budget} give $0<A_f\le2g+2$,
\[
 F_*=S(1-k)/A_f\ge2(1-k),\qquad Q_*=F_*+k/2\ge2-3k/2,
\]
and, keeping the separate coefficients of $u$ and $v$ in the proof of that lemma,
\[
 XOI\ge F_*+\theta u+(1-\theta)v,\qquad
 YOI\ge F_*+(1-\theta)u+\theta v.
\]
With the renaming of sets and coordinates used below, these become
\[
 \widehat XOI\ge F_*+(1-\theta)k+(2\theta-1)t,\qquad
 \widehat YOI\ge F_*+\theta k+(1-2\theta)t.
\]

Suppose first that $\beta>k$. Then the lower bounds on $X$ and $Y$ exceed $k$, by Lemma~\ref{inherit376:actual-cut}, and
$J\ge J_c>\beta>k$. Lemma~\ref{simpleunion:singleton-credit}
requires the lower bound $3/2-k$ on each triple intersection, which is less than $F_*$,
and a lower bound on $XYOI$ which is at most
$\max(2-2k,1+\beta)$. The first is bounded by $Q_*$. To bound the second use
\[
 \bar\beta(k)=\frac{g-1+(n+h+3)k}{2n+h+2g},\quad
 K=\frac{g-1}{n+2g-3}<1/2.
\]
Since $\beta>k$, we have $k<K$, and $\bar\beta(K)=K$.
The affine gap
\[
 \frac{S(1-k)}{2g+2}+\frac{k}{2}-1-\bar\beta(k)
\]
is positive at zero and equals $R_b/[(2g+2)(n+2g-3)]\ge0$ at $K$ \chk{9.K}.
Thus it is positive for $k<K$, and $Q_*\ge1+\beta$.
Hence Lemma~\ref{simpleunion:singleton-credit} gives
$\sum\omega>N/2$.

Suppose now that $\beta\le k$. Let
$(\widehat X,\widehat Y,\widehat O,\widehat I)=(Y,X,I,O)$, with the coordinates $(t,k-t)$ in place of $(u,v)$,
if $B\ge T$, and let these be the original sets, with the same coordinates, if $T>B$.
The lower bounds on $\widehat X$ and $\widehat Y$ are then $L^*>k$ and $L$, and the lower bound
on $\widehat O\widehat I$ is $J$.
Lemma~\ref{inherit376:actual-cut} gives the required comparisons
$L\ge t+j$, $j+a_0\le k-t$, and $J>\beta$; also $L>k/2$ implies $a_0\le k/2$.
These are the assumptions of
Lemma~\ref{top372:multilayer} on single sets and pairs, and its cardinality assumptions follow from
Lemma~\ref{compact390:cuts}. Its required lower bounds on the cyclic triple intersections are at most
$3/2-k$, and that on $XYOI$ is at most $2-3k/2$, so they follow
from $F_*$ and $Q_*$. The differences between the available and the required bounds on the other two triple intersections are
\[
 F_*-3/2+(2-\theta)k+2\theta t+j-a_0,\qquad
 F_*-3/2+(2+\theta)k-2\theta t-a_0.
\]
Both dominate $F_*-3/2+k\ge1/2-k>0$.
Thus all hypotheses of Lemma~\ref{top372:multilayer} hold, which gives
$\sum\omega>N/2$. In both cases, this contradicts $\sum\omega<N/2$, and
Lemma~\ref{general:LP} completes the proof.

\emph{Part~\ref{compact390:bulk-finite:2}.} If $B+T=4g+1$, then
$S=5g+2\ge4(g+1)$ since $g\ge2$. The quantity $R_b$ in part~\ref{compact390:bulk-finite:1} is $3g(n-2)+3>0$ \chk{9.K}. It increases with $S$, because $n+g-2>0$.
Thus the conditions in part~\ref{compact390:bulk-finite:1} hold whenever $B+T\ge4g+1$.
\end{proof}

\subsubsection{Descent bounds for the cyclic intersections}
\label{short392:cyclic-section}

Better lower bounds on the cyclic triple intersections come from mapping $H$ to a path of length four with one shortcut, which is used by one chosen descent in each return. We keep the notation of Section~\ref{compact390:maps-section}, and write
$H_0=B+T$ for the sum of the outer crossing counts.

\begin{lemma}
\label{descent374:full-maps}
\label{descent378:uniform-cyclic}
\label{descent378:first-bottom-cyclic}
Let $j\in\{1,2,3\}$, and in each of $U$ and $V$ choose a step from height $j$ to height $j-1$.
For $i\in\{0,1,2\}$, let $r_i$ be the total number of steps from height $i+1$ to height $i$
before the chosen steps in $U$ and $V$. Then there is a homomorphism from $H$ to the path
$p_0\to p_1\to p_2\to p_3\to p_4$ with the shortcut
$p_{j-1}\to p_{j+1}$, with fibre sizes
\[
 (B-r_0,B+g-r_1,
 g+T+r_0-r_2+2(1-\ind_{j\ge2}),
 T+r_1+2(\ind_{j\ge2}-\ind_{j\ge3}),
 r_2+2\ind_{j\ge3}).
\]
In particular, choosing the last step from one to zero ($j=1$) in both returns, or the last step from three to two ($j=3$) in both returns, gives
\[
 F_1=(2,B,B+g,T+g,T),\qquad
 F_3=(B-r,B+g-m,g+r+2,T+m,T),
\]
for some $0\le r\le B-2$ and $0\le m\le g-2$, which depend on the order of the steps of $U$ and $V$.
Consequently, we obtain homomorphisms from $H$ to $\{x,y,z\}$ in the following five ways. With $F_1$, we may send $(p_0,\ldots,p_4)$ to $(x,y,z,x,y)$ if $z\in XYO$, to $(x,z,y,z,y)$ if $z\in YOI$, and to $(z,x,y,z,x)$ if $z\in XYI$; with $F_3$, we may send it to $(x,y,z,x,y)$ if $z\in XYI$, and to $(x,z,x,z,y)$ if $z\in XOI$. In these five cases, the numbers of positions sent to $x$, $y$ and $z$ are, respectively,
\[
\begin{gathered}
 (T+g+2,\,H_0,\,B+g),\qquad (2,\,H_0+g,\,H_0+g),\qquad (H_0,\,B+g,\,T+g+2),\\
 (H_0+m-r,\,H_0+g-m,\,g+r+2)\qquad\text{and}\qquad (B+g+2,\,T,\,H_0+g).
\end{gathered}
\]
As each homomorphism costs at least $S$, numbers $(A,C,D)$ give the lower bound $(S-Au-Cv)/D$ on the corresponding intersection.
\end{lemma}
\begin{proof}
Consider one of the returns, say $U$, and let $s_i$ be the number of its steps from $i+1$ to $i$ before the chosen step.
The part of $U$ up to the chosen step has
$(1+s_0,1+s_0+s_1,s_1+s_2+\ind_{j\ge2},s_2+\ind_{j\ge3})$ positions at heights $0,1,2,3$.
Send the separating edge before $U$ to $p_0p_1$, the positions of $U$ at height $i$ before the chosen step to
$p_{i+1}$, and those after it to $p_i$. The chosen step then
uses exactly the shortcut $p_{j-1}\to p_{j+1}$, and $U$
ends at $p_0$. Subtracting the counts above from the counts
$(a+1,a+b,b+c,c)$ for the whole of $U$ gives the numbers of positions after the chosen step,
\[
 (a-s_0,a+b-s_1,b+c+s_0-s_2+1-\ind_{j\ge2},
 c+s_1+\ind_{j\ge2}-\ind_{j\ge3},s_2+\ind_{j\ge3}).
\]
and adding the same counts for $V$ gives the fibre sizes, counting both
separating edges. For the last steps from one to zero, and from three to two, we have
$(r_0,r_1,r_2)=(B-2,g,T)$ and $(r,m,T-2)$, respectively.
In each of the five cases, the edges of the path and its
shortcut are present by the stated membership of $z$,
and adding the fibre sizes at each vertex gives the listed multiplicities. Each homomorphism
costs at least $S$, which gives the lower bounds.
\end{proof}

The bound obtained from $F_3$ depends on the unknown numbers $r$ and $m$, and we now remove this dependence.

\begin{lemma}
\label{descent374:full-maps:rest}
For any $\rho$ with $r\le\rho\le B-2$, the homomorphism for $XYI$ obtained from $F_3$ gives the lower bounds
\[
 Q_{I,\rho}=\frac{S-(H_0+g-2-\rho)u-(H_0+2)v}{g+\rho+2}
 \quad(u\ge v),
\]
\[
 Q_{L,\rho}=\frac{S-(H_0-\rho)u-(H_0+g)v}{g+\rho+2}
 \quad(v\ge u).
\]
Choosing instead the first step from one to zero in each return gives
a homomorphism for $XYO$ with multiplicities
$(H_0+m,B+g-m+t,T+g+2-t)$, where
$0\le m\le g$, $0\le t\le T$. Consequently
\[
 P_U=\frac{S-(H_0+g)u-Bv}{T+g+2}\quad(u\ge v),\qquad
 P_V=\frac{S-H_0u-(B+g)v}{T+g+2}\quad(v\ge u)
\]
are lower bounds on $XYO$.
\end{lemma}
\begin{proof}
For the homomorphism from $F_3$, we take the worst value of $m$, namely
$m=g-2$ if $u\ge v$, and $m=0$ if $v\ge u$; the resulting bound
decreases in $r$. For the first steps from one to zero, we have $r_0=0$, $r_1=m$ and $r_2=t$ for some $m$ and $t$,
and we take $m=g$ or $m=0$ in the two orders of $u$ and $v$; the resulting
bound increases in $t$. In each case the sign of the derivative is that
of a linear form in $u,v$ whose coefficients are at most $2S$, which is
positive for $k<1/2$ \chk{9.L}. This gives the bounds with $\rho$, and
the value at $t=0$ gives $P_U,P_V$.
\end{proof}

\begin{lemma}
\label{descent378:top-two-cyclic}
\label{descent378:universal-X-cyclic}
\label{descent378:universal-Y-cyclic}
For all integers $B,T,g\ge2$, the
lower bounds
\[
 F_X=\frac{g+1+(H_0-g-1)u+(2B-1)v}{D_X},\quad D_X=2B+g+2,
\]
\[
 F_Y=\frac{g+1+(2T-3)u+(H_0-g-1)v}{D_Y},\quad D_Y=2T+g
\]
satisfy
\[
 XYO\ge1-u+(1/2-F_X)_+,\qquad
 XYI\ge1-v+(1/2-F_Y)_+.
\]
These bounds remain valid if $F_X$ or $F_Y$ is replaced by a larger lower bound
on $X$ or $Y$, respectively.
\end{lemma}
\begin{proof}
As $XYO\subseteq XO$ and $XYI\subseteq YI$, we already have $XYO\ge1-u$ and $XYI\ge1-v$, so it suffices to prove
$XYO\ge3/2-u-F_X$ and $XYI\ge3/2-v-F_Y$ on the whole
triangle $u,v\ge0$, $u+v\le1/2$.
We use the homomorphisms of Lemmas~\ref{descent374:full-maps} and~\ref{descent374:full-maps:rest}. Let
$C_B=B+g$, $C_T=T+g+2$, $q=B-T-2$ and
\[
 Q_O=[S-H_0k+(B-g-2)u]/C_B,\qquad
 Q_B=[S-H_0u-(B+g)v]/C_T.
\]
For $XYO$, we use the bound $Q_O$ from the last steps from one to zero if $q\le0$ or $g\ge2B-7$, and
otherwise the bounds $P_U$ for $u\ge v$ and $P_V$ for $v\ge u$ from the first steps from one to zero.
For $XYI$, we use $Q_B$ if $B\ge T+2$ or $g\ge2T+2$, and otherwise the
bounds $Q_{I,B-2}$ for $u\ge v$ and $Q_{L,B-2}$ for $v\ge u$.
All bounds are affine in $(u,v)$, so each comparison is checked at the
vertices of the triangle or of its halves $u\ge v$, $v\ge u$; in every
case the vertex values are positive \chk{9.M}.
Finally, replacing $F_X$ or $F_Y$ by a larger lower bound only decreases the right-hand sides.
\end{proof}

In the following variant of Lemma~\ref{short392:incidence-table-two}\ref{short392:incidence-table:R}, the lower bound on $OI$ is traded against the lower bounds on $XO$ and $YI$.

\begin{lemma}
\label{surplus380:coupled}
Let $\omega\ge0$ be a weighting of a set of $N$ vertices, and let $X,Y,O,I$ be sets of at least $d>N/2$ vertices. Let $u,v\ge0$, $k=u+v<1/2$,
$s=1/2-k$ and $q=1-k$. Suppose that $X,Y,O,I$ have lower bounds
$l_X,l_Y,v,u$, respectively, with $l_X,l_Y\ge k$,
that $XY\ge1-\beta$, where $0\le\beta<1/2$,
and that $OI\ge J$. Let
\[
 a_s=(1/2-l_X)_+,\quad b_s=(1/2-l_Y)_+,\qquad
 \tau=\max\{0,a_s+b_s-s,2a_s-u-s,2b_s-v-s\}.
\]
Suppose that
\[
 J\ge\max\{k,\beta,v+a_s-b_s,u+b_s-a_s\},
\]
\[
 XO\ge1-u+\tau,\quad YI\ge1-v+\tau,
\]
\[
 XYO\ge1-u+a_s,\quad XYI\ge1-v+b_s,\quad
                                      XYOI\ge1+s+\tau.
\]
Let $Z_{XI}=\{z\in XI:\omega(z)<1/2\}$ and
$Z_{YO}=\{z\in YO:\omega(z)<1/2\}$, and suppose that the sets
\[
 E_{XI}=Z_{XI}\cup\{z\in O\cup Y:\omega(z)<q\}\qquad\text{and}\qquad
 E_{YO}=Z_{YO}\cup\{z\in I\cup X:\omega(z)<q\}
\]
have at most $N-d$ vertices whenever $Z_{XI}$ or $Z_{YO}$, respectively,
is nonempty.
Then $\sum\omega>N/2$.
\end{lemma}
\begin{proof}
Let $a=a_s$ if $Z_{XI}$ is nonempty and $a=0$ otherwise, let $b=b_s$ if $Z_{YO}$ is nonempty and $b=0$ otherwise,
and let $\lambda=a+b$. Thus $0\le a\le a_s\le s$,
$0\le b\le b_s\le s$, and $0\le\tau\le s$.
Set
\[
 \Delta=\max(\lambda,a_s-u+a,b_s-v+b,\beta-k,a_s+b_s-J),
 \qquad D=\Delta-\lambda\ge0.
\]
The definition of $\tau$ bounds each entry by $s+\tau$, and
the assumptions on $J$, together with $2J\ge k$,
bound each by $J+a_s+b_s-k$. Hence $\Delta\le\min(s+\tau,J+a_s+b_s-k)$. We now choose the coefficients $c_I=i$ and $c_O=o$ in the rectangle
\[
 \max(a_s,u+\Delta-a_s)\le i\le u+\Delta-a,\qquad
 \max(b_s,v+\Delta-b_s)\le o\le v+\Delta-b
\]
whose two intervals are nonempty by the definition of $\Delta$ and
$a\le a_s$, $b\le b_s$. The minimum of $i+o$ over the rectangle is $\max(a_s+b_s,k+2\Delta-a_s-b_s, \Delta+v+a_s-b_s,\Delta+u+b_s-a_s)\le\Delta+J$. where the bound for the first term uses $\Delta\ge a_s+b_s-J$, that for the second the upper bound on $\Delta$, and those for the last two the assumptions on $J$.
The maximum of $i+o$ over the rectangle is $k+2\Delta-\lambda$, which is at least
$\Delta+\max(k,\beta-\lambda)$, by
$\Delta\ge\lambda,\beta-k$. Since $J\ge k,\beta$,
we may choose a point of the rectangle with $\Delta+\max(k,\beta-\lambda)\le i+o\le\Delta+J$. Set $e_A=v+\Delta-b-o$, $e_B=u+\Delta-a-i$, $E=e_A+e_B$.
These choices give
\[
\begin{gathered}
 e_A,e_B\ge0,\quad E\le D,\quad c_I=i\ge a_s,\quad c_O=o\ge b_s,\\
 D+\lambda-E\ge\beta-k,\quad k+D-E\le J,\\
                  a+e_B\le a_s,\qquad b+e_A\le b_s.
\end{gathered}
\]
We now apply Lemma~\ref{short392:incidence-table} with $\mu=0$ and
$\Gamma=\tau$. Its condition $D+\Lambda-E\le2s$ follows from
$\Delta\le s+\tau\le2s$, and its condition $q_0-C(q_0)\ge E-D$ from
$E\le D$. The assumed bounds on $XYO$ and $XYI$ give its conditions on the cyclic triple intersections, as $a+e_B\le a_s$ and $b+e_A\le b_s$, and the bound on $XYOI$ follows from
$\mathcal B=\Delta\le s+\tau$. Thus Lemma~\ref{short392:incidence-table} gives $\sum\omega>N/2$.
\end{proof}
\subsubsection{The case of a large middle crossing count}
Suppose now that $g\ge|B-T|$, that is, the middle crossing count is at least the difference of the outer ones. Then Lemma~\ref{surplus380:coupled} applies, as follows; here $n$ and $n+h$ are the outer crossing counts, in the order which makes the formulas below the bounds of Lemma~\ref{compact390:normalized-cover}.

\begin{lemma}
\label{gap394:canonical-inputs}
Let $n\ge2$, $h\ge0$, $g\ge\max(2,h)$ be integers, and suppose either
$h\ge2n$ or $2n+h\le3g+2$. Put
\[
 a=h+1,\quad D_X=2n+2h+g+2,\quad D_Y=2n+g,
 \quad D=2n+h+2g,\quad S=2n+h+g+1.
\]
Let $\omega\ge0$ be a weighting of a set of $N$ vertices, and let $X,Y,O,I$ be sets of at least $d>N/2$ vertices. Let
$u,v\ge0$, $k=u+v<1/2$ and $s=1/2-k$, and define
\[
\begin{split}
 F_X&=[g+1+(2n+h-g-1)u+(2n+2h-1)v]/D_X,\\
 F_Y&=[g+1+(2n-3)u+(2n+h-g-1)v]/D_Y,\\
 J_c&=[g+1+(2n-3)u+(2n+2h-1)v]/D.
\end{split}
\]
Suppose that $X,Y,O,I$ have lower bounds
$l_X=\max(F_X,k)$, $F_Y$, $v$ and $u$, and that $OI$ and $XY$ have lower bounds
$J$ and $1-\beta$, where $J\ge\max(J_c,k)$,
$0\le\beta<1/2$ and $\beta\le J_c$.
Suppose also that $XO\ge1-u+(1-k)/6$, $YI\ge1-v+(1-k)/6$, that the two cyclic triple
intersections satisfy
\[
 XYO\ge1-u+(1/2-F_X)_+,\qquad
 XYI\ge1-v+(1/2-F_Y)_+,
\]
and that $XYOI$ has both lower bounds
\[
             1+s,\qquad S(1-k)/(2g+2)+k/2.
\]
Finally, suppose that the cardinality bounds of Lemma~\ref{surplus380:coupled} hold. Then
$\sum\omega>N/2$.
\end{lemma}
\begin{proof}
Since $g\ge h$, either assumed inequality gives $2n+h\le3g+2$.
Let
$a_s=(1/2-l_X)_+$ and $b_s=(1/2-F_Y)_+$, and put
\[
 G_X=J_c+F_X-F_Y-v,\qquad G_Y=J_c+F_Y-F_X-u.
\]
We verify the hypotheses of Lemma~\ref{surplus380:coupled},
taking $\tau=\max(0,a_s+b_s-s,2a_s-u-s,2b_s-v-s)$.
All comparisons below are affine in $(u,v)$ and are checked at the
vertices of the stated polygons \chk{9.N}.

We first prove the required bounds on $J$. Both lower bounds $l_X$ and $F_Y$
are at least $k$, since $F_Y-k\ge0$ at the vertices of the triangle; so
the assumed bounds on the cyclic triple intersections are those required in Lemma~\ref{surplus380:coupled}.
Moreover $G_Y>0$ throughout the triangle. If $b_s>0$, then
$b=2n-g-2\ge0$, and the region $F_Y\le1/2$ is then the triangle with
vertices $O=(0,0)$, $U=(b/(2c_u),0)$, $V=(0,b/(2c_v))$, where
$c_u=2n-3>0$ and $c_v=b+h+1>0$; on it $G_X>0$.
These two comparisons give the required bound on $J$
when $a_s$ and $b_s$ are both positive. If instead $b_s>a_s$, then the
maximum with $k$ cannot replace $F_X$, since $F_Y>k$, so the
comparison $G_Y>0$ still applies.

The only remaining nontrivial target has
$a_s>0=b_s$ and $v+a_s>k$, with polygon $F_Y\ge1/2$, $F_X+u<1/2$.
There $u<1/4$, since $F_X+u-1/2=(2g+1-h)/(4D_X)>0$ at $(1/4,0)$, and so
$D(J_c-k)=g+1-(2g-h+1)k-2au>a(1/2-2u)>0$. The affine margin
$G=J_c+F_X-v-1/2$ equals $G_X>0$ on $F_Y=1/2$, and is positive on
$F_X+u=1/2$ by $J_c>k$. The polygon avoids $k=1/2$, since $F_X+u>1/2$
at both half axes; if it contains the origin, then $g\ge2n-2$ and
$G(O)>0$. Thus $G>0$ at every vertex. The other required lower bounds
on $J$ are at most $u,v$, so they follow from $J\ge k$.

To bound $\tau$, put $q=2n+2h-3g-2$ and $\tau_0=q_+/(2D_X)$; the assumed
inequality gives $q\le h$ and $\tau_0<1/6$. We have
\[
 1/2+v-2F_X=[qs-D_Yu-2(g-1)v]/D_X\le\tau_0(1-2k),
 \qquad 1/2+u-2F_Y\le0.
\]
If $a_s$ or $b_s$ vanishes, then the corresponding
terms in the definition of $\tau$ are nonpositive.
When both are positive, $a_s+b_s-s$ is equal to the affine function
$1/2+k-F_X-F_Y$ on the triangle $O,U,V$. It is at most $\tau_0$ at $O$,
and at $U,V$ it equals $k-F_X\le0$.
Thus $\tau\le\tau_0(1-2k)$, and since $\tau_0<1/6$, the bounds on
$XO,YI$ suffice.

For $q\le0$ we use the first bound on $XYOI$.
For $q>0$, the other lower bound is $c(1-k)+k/2$
with $c=S/(2g+2)$, and $c>3/2+\tau_0$ \chk{9.N}.
Since $c(1-k)+k/2-(1+s)=(c-3/2)(1-k)$,
the second bound on $XYOI$ is at least $1+s+\tau$.
Lemma~\ref{surplus380:coupled} now gives $\sum\omega>N/2$.
\end{proof}

When $X\ge1/2$ but the lower bound on $OI$ is less than $k$, we use instead the following special case of Lemma~\ref{short392:incidence-table}.

\begin{lemma}
\label{gap390:regular-jbad-count}
Let $\omega\ge0$ be a weighting of a set of $N$ vertices, and let $X,Y,O,I$ be sets of at least $d>N/2$ vertices. Let $u,v\ge0$, $k=u+v<1/2$, $s=1/2-k$ and
$h_c=1-k$, and suppose that
\[
 X\ge1/2,\quad Y\ge l_Y\ge k,\quad O\ge v,\quad I\ge u,
 \qquad OI\ge J,\quad \max(u,v)\le J<k.
\]
Let $b_s=(1/2-l_Y)_+$ and $e=k-J$, and suppose that
\[
 b_s\le v,\qquad 0\le\beta<J,\qquad
 XY\ge1-\beta,\quad XO\ge1-u,\quad YI\ge1-v,
\]
\[
 XYO\ge1-u+e,\qquad XYI\ge1-v+b_s,\qquad
 XYOI\ge1+s.
\]
Let $Z_{YO}=\{z\in YO:\omega(z)<1/2\}$ and
$E_{YO}=Z_{YO}\cup\{z\in I\cup X:\omega(z)<h_c\}$, and suppose that
$|E_{YO}|\le N-d$ if $Z_{YO}$ is nonempty.
Then $\sum\omega>N/2$.
\end{lemma}
\begin{proof}
In Lemma~\ref{short392:incidence-table}, take
$a=D=e_A=\mu=0$ and $e_B=e$, and let $b=b_s$ if $Z_{YO}$ is nonempty and $b=0$ otherwise; note that $Z_{XI}$ is empty as $X\ge1/2$.
Here $a_s=0$, $b_s\le s$ and $e\le\min(u,v)$, so
$c_I=u+b-e\ge0$ and $c_O=v\ge b_s$.
The conditions on $a$ and $b$ hold.
We also have $b-e\ge\beta-k$, $b-e\le2s$ and $k-e=J$,
while $q_0\ge k-\beta>e$ gives the required threshold inequality.
Take $\Gamma=0$; the pair inequalities then follow from $\mathcal B=b\le s$.
The bounds on $XYO,XYI$ imply the required additions $e$ and
$b\le b_s$, and the bound on $XYOI$ follows from $1+\mathcal B\le1+s$.
Thus Lemma~\ref{short392:incidence-table} applies.
\end{proof}

The conditions of Lemma~\ref{gap394:canonical-inputs} hold when $B\ge T$, as follows.

\begin{lemma}
\label{gap390:bottom-preliminaries}
\label{gap392:bottom-inputs}\label{allbalance377:actual-inputs}
Let $B=n+h$ and $T=n\ge2$, where $h\ge0$ and $g\ge\max(2,h)$, and
suppose that
$m=R+\iota\ge h$. Write
\[
 a=h+1,\quad D_X=2n+2h+g+2,\quad D_Y=2n+g,
 \quad D=2n+h+2g,\quad S=2n+h+g+1.
\]
Then the cover $w$ satisfies the following.
\begin{props}
\item\label{gap390:bottom-preliminaries:1} The lower bounds of Lemma~\ref{compact390:normalized-cover}\ref{compact390:normalized-cover:2} are
\[
\begin{split}
 F_X&=\frac{g+1+(2n+h-g-1)u+(2n+2h-1)v}{D_X},\\
 F_Y&=\frac{g+1+(2n-3)u+(2n+h-g-1)v}{D_Y},\\
 J_c&=\frac{g+1+(2n-3)u+(2n+2h-1)v}{D},\qquad
 \beta=\frac{g-1+(n+1)u+(n+h+3)v}{D}.
\end{split}
\]
Moreover, $l_X=\max(F_X,k)$ and $l_Y=F_Y$
are lower bounds on $X$ and $Y$, there is a lower bound
$J>\max(k,\beta)$ on $OI$, and $XO$ and $YI$ have lower bounds
$1-u+(1-k)/6$ and $1-v+(1-k)/6$, respectively.
\item\label{gap390:bottom-preliminaries:2} If $h\ge2n$ or $2n+h\le3g+2$, then all hypotheses of
Lemma~\ref{gap394:canonical-inputs} hold, and so the cover $w$ does not exist.
\end{props}
\end{lemma}
\begin{proof}
\emph{Part~\ref{gap390:bottom-preliminaries:1}.} Part~\ref{compact390:normalized-cover:2} of Lemma~\ref{compact390:normalized-cover} gives $0<\beta<1/2$ and
$D(J_c-\beta)=2+(n-4)u+(n+h-4)v>0$.
At the three vertices of the triangle $u,v\ge0$, $u+v\le1/2$, $F_Y-k$ is positive
\chk{9.N}, and the lower bound $F_R$ on $X$ exceeds $k$ by
\[
 g+1+R-\iota,\qquad (m-h+1)/2,\qquad (g-1+R-\iota)/2
\]
over its denominator \chk{9.O}. These are positive, since $\iota\le g-2$ and $m\ge h$,
so $l_X$ is a lower bound on $X$.
For $z=m-h+1\ge1$, combining the two lower bounds $J_c$ and $J_R$ on $OI$ with weights $z$ and $a$ gives the lower bound
\[
 J_w=k+
 \frac{[z(g+1)+a(g+1+R-\iota)](1-2k)}
 {zD+a(D+R-\iota)}>k,
\]
since their $u$ coefficients cancel \chk{9.O}.
Take $J=\max(J_c,J_w)$.

The positions at heights zero and three of the returns give at least $B$ sources and at least $T$ sinks, and, similarly, $H$ has at least
$g$ sinks; moreover, $H$ has as many sources as sinks.
Choosing these sources or sinks in the larger bipartition class in the construction of
Lemma~\ref{compact390:actual-reserves}, with $M_*=\max(B,T,g)$,
gives the additional terms $M_*(1-k)/(2S-M_*)$ in the bounds on $XO$ and $YI$.
Since $M_*\ge(S-1)/3$ and $S\ge7$, we have $7M_*\ge2S$, and so these terms are at least $(1-k)/6$.

\emph{Part~\ref{gap390:bottom-preliminaries:2}.} Part~\ref{gap390:bottom-preliminaries:1} gives the bounds involving
$\max(F_X,k)$, $J\ge\max(J_c,k)$ and $\beta_c<J_c$.
The lower bounds on the two cyclic triple intersections are given by
Lemma~\ref{descent378:top-two-cyclic}. The bound~\eqref{balanced:bound}
and Corollary~\ref{geom368:canonical-fold-budget} give both lower bounds on $XYOI$, and the cardinality bounds are given by
Lemma~\ref{compact390:cuts}. Thus Lemma~\ref{gap394:canonical-inputs} applies.
\end{proof}

For the case $T>B$, we need a technical lemma, which reduces an inequality of the form $\max_iA_i>\Psi$ on a region of the $(k,u)$-plane to the same inequality on a curve.

\begin{lemma}
\label{gap396:affine-cap-witness}
Let $a,D,E,A_0,c>0$, $Z=D+2a=E+A_0$ and $L=2c+a$.
For $c/L<k<1/2$, put $w=Lk-c$, $s=1/2-k$, $u_0=w/Z$ and
\[
 e=\min(u,(w-2au)/D),\qquad t=(w-A_0u)/E,
 \qquad 0\le u<\min(w/(2a),w/A_0).
\]
Let $\Psi$ be either $e+t$ or $e+t\min(1,e/s)$,
and let $A_i(k,u)$ be finitely many affine functions. For each $\xi\ge0$,
consider the point with $u=u_0$ given by
\[
 k_\xi=\frac{c+Z\xi/2}{L+Z\xi},\quad
 u_\xi=\frac{a\xi}{2(L+Z\xi)},\quad e=t=u_\xi,
\]
and suppose that some $i$ satisfies
\[
 A_i(k_\xi,u_\xi)\ge\Psi(k_\xi,u_\xi),\qquad
 A_i(1/2,0)\ge a/(2E),
\]
with the first inequality strict for $\xi>0$.
If $\Psi=e+t$, suppose also that at least one of these inequalities is strict at $\xi=0$.
Then $\max_iA_i>\Psi$ whenever $u\le u_0$.
Moreover, let $u_*=\min(w/(2a),w/A_0)$.
Then $\max_iA_i>\Psi$ throughout the region above if each $A_i$ used above has
$A_i-e-t$ nondecreasing on $[u_0,u_*]$ and, if $\Psi=e+t\min(1,e/s)$, also $A_i(k,u_*)\ge e(k,u_*)$.
\end{lemma}
\begin{proof}
For $u\le u_0$, let $\xi=u/s$, so that $k\ge k_\xi$ and $e=u$. The points with this value of $\xi$ form a segment from $(k_\xi,u_\xi)$ to $(1/2,0)$, on which
both choices of $\Psi$ are affine:
the second is $\xi(s+t)$ for $\xi\le1$, and $e+t$ otherwise.
Their limits at $(1/2,0)$ are at most $a/(2E)$, so the inequalities at the two ends of the segment
give $A_i>\Psi$ on it. At $\xi=0$, $k>k_0$;
the required strict endpoint for $e+t$ suffices, while for the
other right-hand side the limiting value is zero and
$A_i(1/2,0)\ge a/(2E)>0$.
For $u\ge u_0$, fix $k$; at $u=u_0$ we have
$\xi=w/(Zs)>0$. Monotonicity proves the assertion for $\Psi=e+t$,
and for the second choice of $\Psi$ while $e\ge s$.
On the remaining interval, $\Psi=e+et/s$ is convex, since
$\Psi''=2e't'/s>0$. Thus $A_i-\Psi$ is concave, positive at its
initial point and nonnegative at $u_*$, where $et=0$ and
$\Psi=e$. The endpoint hypothesis therefore gives strict positivity before $u_*$.
\end{proof}

\begin{lemma}
\label{gap394:top-inputs}
Let $B=n\ge2$ and $T=n+h$, where $h\ge1$, $g\ge\max(2,h)$ and
$2n+h\le3g+2$. Then the cover $w$ does not exist.
\end{lemma}
\begin{proof}
We use the lower bounds of
Lemma~\ref{compact390:normalized-cover}\ref{compact390:normalized-cover:2}, and let
\[
 a=h-1,\quad D_X=2n+g+2,\quad E=2n+2h+g,
 \quad D=2n+h+2g,\quad L=2g+h+1.
\]
Always $0<\beta_c<1/2$, $J_c>\beta_c$ and $F_X\ge k$ \chk{9.P}, and
\[
 D(J_c-k)=g+1-Lk+2au,\qquad
 E(F_Y-k)=g+1-Lk+(g+h-2)u.
\]
Since $g+h-2\ge2a$, the inequality $J_c\ge k$ implies $F_Y\ge k$.

Suppose first that $J_c\ge k$ and $h\ge2$. We reduce to Lemma~\ref{gap394:canonical-inputs}, after renaming the sets and coordinates by
\[
 (X',Y',O',I')=(Y,X,I,O),\qquad (u',v')=(v,u),
 \quad n'=n+1,\quad h'=h-2.
\]
The functions $F_X,F_Y,J_c$ in
Lemma~\ref{gap394:canonical-inputs} are then exactly
$F_Y,F_X,J_c$, with the same $S$ and the same inequality
$2n'+h'=2n+h\le3g+2$. Indeed, the formulas for $F_X$ and $F_Y$ are those for outer crossing counts $B+1$ and $T$, and exchanging these exchanges the two formulas. We keep the bound $\beta_c\le J_c$.
The lower bounds on the cyclic triple intersections and the cardinality bounds are
exchanged together with the sets. The additional terms in the bounds on $XO$ and $YI$ are at least
$(1-k)/6$: as in the proof of Lemma~\ref{gap390:bottom-preliminaries}, $H$ has at least
$M_*=\max(B,T,g)\ge(S-1)/3$ sources and as many sinks, and the
construction of Lemma~\ref{compact390:actual-reserves} gives $M_*(1-k)/(2S-M_*)\ge(1-k)/6$.
The bound~\eqref{balanced:bound} and Corollary~\ref{geom368:canonical-fold-budget} give the two required
lower bounds on $XYOI$, so Lemma~\ref{gap394:canonical-inputs} applies.

If $h=1$, then $J_c-k=(g+1)(1-2k)/D>0$, $D_X=E$ and
$F_X-F_Y=(g-1)(v-u)/D_X$. The differences $J_c+F_X-F_Y-v$ and
$J_c+F_Y-F_X-u$, and the inequalities $2F_X\ge1/2+v$, $2F_Y\ge1/2+u$
and $F_X+F_Y\ge1/2+k$, are affine and hold at the three vertices.
Together with $F_X,F_Y\ge k$, these give $\tau=0$, and the lower bounds on the
cyclic triple intersections and on $XYOI$ imply the hypotheses of
Lemma~\ref{surplus380:coupled} directly.

It remains to consider $J_c<k$, which forces $h\ge2$.
Put $Z=D+2a$ and $A_0=g+h-2$, and
\[
 \ell=Lk-(g+1)\in(0,a/2),\quad
 J=\max(J_c,u,v),\quad e=k-J=\min\{u,(\ell-2au)/D\},
 \quad b_s=(1/2-F_Y)_+.
\]
The region where $J_c\le k$ is the triangle with vertices
$A=(0,(g+1)/L)$, $V=(0,1/2)$ and $C=(1/4,1/4)$.
At these vertices $F_X-1/2$ is $N_X/(2LD_X)$, $(g-1)/(2D_X)$ and
$(g+h-2)/(4D_X)$, where $N_X=2g^2+(h-1)(g-2n)-2>0$ by the assumed
inequality \chk{9.P}; so $F_X\ge1/2$ there.
We also have $J_c>u$ and $u<\ell/(2a)<k/2$, while
$k>(g+1)/L>1/3$. The values of $b_s$ at $A,V,C$ are
$a/(2L),a/(2E),0$, all below $1/6$,
so $b_s<v$ throughout the region.

The lower bounds on $XO$ and $YI$ have the additional term
\[
 \Gamma=\frac{M_*}{2S-M_*}(1-k),\qquad M_*=\max(T,g),
\]
and $\Gamma>b_s$ on this polygon.
Indeed, $H$ has at least $T$ sinks at height three, and similarly at least $g$ sinks, and as many sources as sinks, and we use the larger bipartition class in
Lemma~\ref{compact390:actual-reserves}. Using $M_*\ge T$ at $A$,
$M_*\ge(T+g)/2$ at $V$, and $b_s=0$ at $C$, the comparison holds at the
vertices \chk{9.P}, and affineness proves it.
The same bound gives $M_*/(2S-M_*)>a/E$.
Since $e\le\ell/Z<a(1-k)/Z$ and $Z=E+A_0>E$, it also
gives $\Gamma>e$. Therefore the required bounds on both cyclic triple intersections follow from those on $XO$ and $YI$.

For $F_Y\ge k$ we have $u\ge\ell/A_0>0$, so $J<k$;
also $b_s\le s$.
The bounds on $XO$ and $YI$ give $XYO\ge1-u+e$ and
$XYI\ge1-v+b_s$, and~\eqref{balanced:bound} gives $XYOI\ge1+s$.
These are the hypotheses of
Lemma~\ref{gap390:regular-jbad-count}.

Finally suppose that $F_Y<k$. Then $\beta_c<J_c<k$, $F_Y>u$,
and $t=k-F_Y=(\ell-A_0u)/E>0$.
The remaining region is the triangle with vertices $A,V$ and
$W=(a/(2A_0),1/2-a/(2A_0))$.
The lower bound $Q_O=[S-(2n+h)k+(n-g-2)u]/(n+g)$ on $XYO$ from Lemma~\ref{compact390:actual-reserves} satisfies
$Q_O+F_Y>2-k$ at these three vertices \chk{9.P}, hence throughout.
Thus $XYO$ and $XYOI$ have
lower bound $1+s+b_s$, so all hypotheses of
condition~\ref{short392:incidence-table:W} of Lemma~\ref{short392:incidence-table-two} now hold.
It remains to prove its last condition, and we prove the stronger inequality with right-hand side $e+t$,
using Lemma~\ref{gap396:affine-cap-witness}
with $c=g+1$, $a=h-1$, $Z=D+2a=E+A_0$ and $u_0=\ell/Z$;
the region is
$0\le u<\min(\ell/(2a),\ell/A_0)$.
At $u=u_0$, put $\xi=e/s$ and $K=L+Z\xi$. Multiplication by
the indicated positive denominators gives
\[
\begin{split}
 2DK\{k-\beta_c-e-t\}&=P+C\xi,\quad
 P=2ng+4h+2n-4,\\
 C&=2n^2+2ng+(n-2g)h+2g,\\
 2D_XK\{F_X-(1-k)-e-t\}&=2B_F+Q_F\xi,\quad
 B_F=g^2-(2n+1)h+2n,\\
 Q_F&=2g^2+(2n-1)g+h^2-4nh-8h+2n+6.
\end{split}
\]
For $g\le2n$ we use $k-\beta_c$. Then $C>0$, being affine in
$0\le h\le g$ with positive endpoint values; the limiting difference
from $e+t$ at $(k,u)=(1/2,0)$ is positive, and for $u\ge u_0$ the
difference has positive derivative \chk{9.P}.
For $g\ge2n$ we use $F_X-(1-k)$. Then $B_F\ge0$ and $Q_F>0$, the
limiting difference at $(1/2,0)$ is
$H_F/(2D_XE)$ with $H_F=(g-1)E-aD_X>0$, and for $u\ge u_0$ the
derivative is positive \chk{9.P}.
The inequality is strict at $u_0$ for positive $\xi$, and at $\xi=0$
it is nonnegative, with a positive limit at $(1/2,0)$. Lemma~\ref{gap396:affine-cap-witness} also covers the points with $u=e=0$, whichever
of $e$ and $t$ vanishes first.
It gives $\max(k-\beta_c,F_X-(1-k))>e+t\ge e+t\min(1,e/s)$.
Thus condition~\ref{short392:incidence-table:W} of Lemma~\ref{short392:incidence-table-two} holds, and gives $\sum\omega>N/2$. This covers all values of $u$ and $v$, and proves the lemma.
\end{proof}

\begin{theorem}
\label{compact390:gap-finite}
\label{compact390:top-gap-finite}\label{compact390:bottom-gap-finite}
If $g\ge|B-T|$, then every digraph $F$ with $\delta^0(F)>|F|/2$ has a
perfect fractional homomorphic $H$-tiling.
\end{theorem}
\begin{proof}
Let $n=\min(B,T)$ and $h=|B-T|\le g$.
We first treat the case $2n+h>3g+2$, in either order of $B$ and $T$.
Then $g<2n$ and, as these are integers, $S\ge4(g+1)$.
At $S=4(g+1)$ the second expression in
part~\ref{compact390:bulk-finite:1} of Theorem~\ref{compact390:bulk-finite}
is $(g+1)(2n-g-1)\ge0$, and it increases with $S$.
So that theorem applies, and we may assume that $2n+h\le3g+2$.

Suppose first that $B\le T$. If $h=0$, then $R+\iota\ge h$, so part~\ref{gap390:bottom-preliminaries:2} of Lemma~\ref{gap392:bottom-inputs} shows that there is no deficient cover,
and if $h\ge1$, Lemma~\ref{gap394:top-inputs} shows this. Thus Lemma~\ref{general:LP} proves the theorem when $B\le T$.

Now suppose that $B>T$, so that $B=n+h$ and $T=n\ge2$, and let
$m=R+\iota$. If $m\le h$, then, by Lemma~\ref{compact390:literal-converse}, the converse of $H$ has crossing counts
\[
 B'=n+\iota,\quad T'=n+h-R,\quad g'=g+R-\iota\ge2.
\]
so that $B'\le T'$, with $h'=T'-B'=h-m\ge0$ and $g'-h'=g+2R-h\ge0$.
Thus the preceding case applies to the converse of $H$, in the converse of $F$, and Lemma~\ref{compact390:literal-converse} transfers the tiling back to $H$.
If $m>h$, then, as $2n+h\le3g+2$, part~\ref{gap390:bottom-preliminaries:2} of Lemma~\ref{gap392:bottom-inputs}
shows that there is no deficient cover, and Lemma~\ref{general:LP} completes the proof.
\end{proof}


\subsection{Reduction of the remaining crossing counts}
For the crossing counts not covered by Theorems~\ref{compact390:bulk-finite} and~\ref{compact390:gap-finite}, applied to $H$ or to its converse, the top crossing count exceeds the bottom one by more than the middle one, as follows.

\begin{lemma}
\label{compact390:count-reduction}
Suppose that neither $H$ nor its converse satisfies the
conditions of Theorem~\ref{compact390:gap-finite} or of Theorem~\ref{compact390:bulk-finite}.
Then $H$ or its converse, in the form given by Lemma~\ref{compact390:literal-converse}, has $B=n\le T=n+h$ with $h>g\ge2$.
For this orientation, $q=n-R\ge2$ and
\[
 h+5\iota\le3n+4g-5q.
\]
Moreover, if $g<2n$, then $h\le3g-2n+2$,
and if $g\ge2n$, then $h\le4g-2n$.
\end{lemma}
\begin{proof}
If $B\le T$, we use $H$ itself. Otherwise, let $d=B-T>0$; then
$g<d$, as Theorem~\ref{compact390:gap-finite} does not apply. The converse has
$B'=T+\iota$, $T'=B-R$, $g'=g+R-\iota$ and
$g'-(B'-T')=d+g-2\iota\ge d-g+4>0$, so, as the same theorem does not apply to the converse,
$B'<T'$; we use the converse. In either case, $h>g$, as Theorem~\ref{compact390:gap-finite} does not apply.

As the crossing counts are integers, the failure of part~\ref{compact390:bulk-finite:2} of
Theorem~\ref{compact390:bulk-finite} for the converse of the chosen orientation is equivalent to
$h+5(\iota-R)\le4g-2n$ \chk{9.A}, which is the first asserted inequality.
The expression $R_b$ in part~\ref{compact390:bulk-finite:1} increases
with $S=2n+h+g+1$. It equals $(g+1)(2n-g-1)$ at $S=4(g+1)$, and
$g(3n-7)-n+5>0$ at $h=4g-2n$ when $g\ge2n$ and $n\ge3$ \chk{9.A}.
So failure of part~\ref{compact390:bulk-finite:1} gives $h\le3g-2n+2$
if $g<2n$, and $h<4g-2n$ if $g\ge2n$ and $n\ge3$. If $n=2$, then $R\le B-2=0$ gives
$q=2$, and the first inequality gives $h\le4g-4$.
\end{proof}

\subsection{The remaining crossing counts}\label{sec:remaining}
By Lemma~\ref{compact390:count-reduction}, it remains to consider $B=n$ and $T=n+h$ with $h>g$, and, by Proposition~\ref{compact390:source-criterion}, it suffices to verify its inequality in $u$ and $v$, together with its two inequalities in the crossing counts. We first reduce the inequality in $u$ and $v$ to an inequality along a curve, and then to the positivity of two polynomials in one variable.

\subsubsection{The inequality along a curve}
Throughout this subsection, let
\[
 a=h-1,\quad D=2n+h+2g,\quad E=2n+2h+g,\quad
 L=2g+h+1,\quad Z=D+2a,\quad A_0=g+h-2.
\]
We use the numbers $R,\iota$ of the orientation chosen in Lemma~\ref{compact390:count-reduction}, and write
\[
 q=n-R\ge2,\quad \delta=\iota-R,\quad A_*=2n+g,
 \quad A_R=A_*-R-\iota,
\]
\[
 P=2ng+4h+2n-4,\quad C=nD-2ga,
 \quad B_0=g^2-(2n-1)h+2n-2,
\]
\[
 Q_0=h^2-(4n+2)h+(2n-1)g+2g^2+2n,
 \quad d=D-2a,
\]
\[
 B_F=B_0+2aR-(g+1)\delta,\quad
 Q_F=Q_0+6aR-d\delta,\quad
 H_F=\{A_R-(2q+1)\}E-aA_R.
\]

\begin{lemma}
\label{compact392:cap-formulas}
Let $u,v\ge0$ with $k=u+v<1/2$, and suppose that $F_Y<k$ and $\max(J_c,u,v)<k$. Let $s=1/2-k$,
$e=k-\max(J_c,u,v)>0$ and $t=k-F_Y>0$. Then
\[
 \ell=Lk-(g+1)\in(0,a/2),\quad
 e=\min\{u,(\ell-2au)/D\},\quad t=(\ell-A_0u)/E.
\]
The two expressions for $e$ agree at $u_0=\ell/Z$;
this precedes $u_J=\ell/(2a)$, where $e$ vanishes. At $u_0$, with $\xi=e/s$ and $K=L+Z\xi$,
\[
 k=\frac{g+1+Z\xi/2}{K},\qquad
 u=e=t=\frac{a\xi}{2K},\qquad s=\frac a{2K}.
\]
For $0\le\xi\le1$, put $\chi=e+t\xi$. Then
\begin{align}
 2DK(k-\beta_c-\chi)
 &=T_c(\xi):=P+C\xi+Da\xi(1-\xi),
 \label{compact392:triangle-cap}\\
 2A_RK(F_R-(1-k)-\chi)
 &=F(\xi):=2B_F+Q_F\xi+A_Ra\xi(1-\xi).
 \label{compact392:singleton-cap}
\end{align}
Both polynomials are strictly concave in $\xi$.
For $\xi\ge1$, the right-hand side $e+t\min(1,e/s)$ in Proposition~\ref{compact390:source-criterion} is $e+t$, and
using $e+t$ for all $\xi$ replaces these two polynomials by $P+C\xi$ and
$2B_F+Q_F\xi$, respectively.
\end{lemma}
\begin{proof}
Substitute the lower bounds of part~\ref{compact390:normalized-cover:2}
of Lemma~\ref{compact390:normalized-cover}; the inequalities $J_c>u$
and $A_0<2a$ select the stated branches \chk{9.B}.
\end{proof}

\subsubsection{The case $C\ge0$}

\begin{lemma}
\label{compact390:nonnegative-slope}\label{compact390:ratio-counts}
Let $F_Y$, $J_c$ and $\beta_c$ be as in Lemma~\ref{compact390:normalized-cover}, with $B=n$ and $T=n+h$,
and let
\[
 C=2n^2+2ng+(n-2g)h+2g.
\]
\begin{props}
\item\label{compact390:nonnegative-slope:1} Let $n\ge2$ and $h>g\ge2$ be integers.
For $u,v\ge0$ with $k=u+v<1/2$, suppose that
\[
 F_Y<k,\qquad \overline J=\max(J_c,u,v)<k.
\]
Write $s=1/2-k$, $e=k-\overline J$ and $t=k-F_Y$.
If $C\ge0$, then
\[
                         k-\beta_c>e+t.
\]
In particular, $k-\beta_c>e+t\min(1,e/s)$.
\item\label{compact390:nonnegative-slope:2} Let $n\ge2$ and $n\le g<h\le2n$ be integers, and let
\[
 N_1=4g^2+2ng+6g+4n+4-4n^2+h(g-2n),
 \qquad N_X=2g^2+(h-1)(g-2n)-2.
\]
Then $N_1>0$, $N_X\ge0$ and $C>0$. Hence Proposition~\ref{compact390:source-criterion} applies, and every digraph $F$ with $\delta^0(F)>|F|/2$ has a perfect fractional homomorphic $H$-tiling.
\end{props}
\end{lemma}
\begin{proof}
For part~\ref{compact390:nonnegative-slope:1}, at $u=u_0$ the difference from $e+t$ has numerator
$P+C\xi>0$, since $P>0$ and $C\ge0$, and its limit at $(1/2,0)$ is
positive \chk{9.D}. Lemma~\ref{compact390:cap-propagation}, applied with
$k-\beta_c$, gives the assertion.

The three signs in part~\ref{compact390:nonnegative-slope:2} hold
throughout the range \chk{9.C}, and $N_1$ is the quantity of
Proposition~\ref{compact390:source-criterion}. Thus the inequalities in the crossing counts
in that proposition hold, and part~\ref{compact390:nonnegative-slope:1} gives
its inequality in $u$ and $v$.
\end{proof}

\subsubsection{Reduction to two polynomials}

The following two simple tests show that the maximum of two polynomials is positive.

\begin{lemma}
\label{cap394:affine-overlap}
Let $P,Q>0$, $B,C\in\mathbb R$ and $A,D\ge0$. Put
\[
 T(\xi)=P+C\xi+D\xi(1-\xi),\qquad
 G(\xi)=2B+Q\xi+A\xi(1-\xi).
\]
Then each of the following conditions implies that $\max(T,G)>0$ on $[0,1]$
and $\max(P+C\xi,2B+Q\xi)>0$ for $\xi\ge1$.
\begin{props}
\item\label{comparison:linear}
$C\ge0$, or $B\ge0$, or $PQ-2CB>0$.
\item\label{comparison:quadratic}
For some $0<\theta<1$, one has $T(\theta)>0$, $G(\theta)>0$
and $G(1)>0$.
\end{props}
Under~\ref{comparison:linear}, we even have $\max(P+C\xi,2B+Q\xi)>0$ for every $\xi\ge0$.
\end{lemma}
\begin{proof}
For~\ref{comparison:linear}, we may assume that $C<0$ and $B<0$.
Then the two affine functions are equal at a positive value of $\xi$, where their
common value is $(PQ-2CB)/(Q-C)>0$, so their maximum is positive for all $\xi\ge0$.
On $[0,1]$, the quadratic terms are nonnegative.
For \ref{comparison:quadratic}, concavity gives $T>0$ on
$[0,\theta]$ and $G>0$ on $[\theta,1]$; beyond one,
$2B+Q\xi$ increases from $G(1)>0$.
\end{proof}

It suffices to verify the inequality of Proposition~\ref{compact390:source-criterion} at $u=u_0$, as follows.

\begin{lemma}
\label{compact390:cap-propagation}
Use the notation of Lemma~\ref{compact392:cap-formulas},
with integers $n\ge2$ and $h>g\ge2$.
Suppose that, for every $\xi>0$, at $u=u_0=\ell/Z$ one of $k-\beta_c$ and $F_R-(1-k)$
is greater than $e+t\min(1,\xi)$, and that $H_F\ge0$ if $F_R-(1-k)$ is needed for some $\xi$.
Then $\max\{k-\beta_c,F_R-(1-k)\}>e+t\min(1,e/s)$ at every point with
$F_Y<k<1/2$ and $\max(J_c,u,v)<k$.
Moreover, if $e+t\min(1,\xi)$ is replaced by $e+t$ in the hypothesis, then $e+t\min(1,e/s)$ may be replaced by $e+t$ in the conclusion.
\end{lemma}
\begin{proof}
Apply Lemma~\ref{gap396:affine-cap-witness} with $c=g+1$.
At $(k,u)=(1/2,0)$ the two differences from $e+t$ have limits
\[
 \frac{h(h-g-1)+2n(n+h)+g(n+1)}{2DE}>0,
 \qquad \frac{H_F}{2A_RE}\ge0,
\]
and for $u\ge u_0$ both have positive derivative in $u$ \chk{9.D}.
Since $A_0<2a$, the other endpoint has $e=0$ and $J_c=k>\beta_c$, so
the first expression is positive there; the second is
positive by monotonicity of $F_R$ and positivity at $u_0$.
At $\xi=0$, we use $k-\beta_c$, whose numerator is $P>0$.
\end{proof}

\subsubsection{The case $h>2n$}
It remains to check the hypotheses of Lemma~\ref{compact390:cap-propagation} when $h>2n$. First, some of the coefficients have the right signs.

\begin{lemma}
\label{compact390:middle-endpoints}
\label{compact390:high-endpoints}
Let $n\ge2$, $2\le q=n-R\le n$ and
$h>g\ge2$ be integers, and suppose that $h+5\iota\le3n+4g-5q$, and that one of the following holds.
\begin{props}
\item $n\le g<2n$ and $2n<h\le3g-2n+2$.\label{compact390:middle-endpoints:M}
\item $g\ge2n$.\label{compact390:middle-endpoints:H}
\end{props}
Then $H_F\ge0$ and $Q_F>0$, and the first inequality is strict except at $(n,g,h,R,\iota)=(2,3,7,0,0)$.
Moreover, in~\ref{compact390:middle-endpoints:M}, we have $N_1,N_X>0$.
\end{lemma}
\begin{proof}
Let $X=E-a$ and $A=A_R=n+g+q-\iota$. The assumed bound on $\iota$ gives
$5A\ge X+10q-1$, so
\[
 5H_F\ge X^2-6X-(10n+5)a
 =\Bigl(h+g-3n-\frac92\Bigr)^2+15n(n-1)-\frac{81}4+(10n+5)(g-2n).
\]
The right-hand side is positive in~\ref{compact390:middle-endpoints:H},
and in~\ref{compact390:middle-endpoints:M} when $n\ge4$, since there
$g\ge(4n-1)/3$. The cases of~\ref{compact390:middle-endpoints:M} with
$n\le3$ are finitely many and are checked directly \chk{9.E}.

For $Q_F$, write $Q_F=Q_0(q)-d\iota$, where
$Q_0(q)=Q_F|_{\iota=0}$ decreases in $q$. Bounding $\iota$ below by zero if
$d\le0$, and above by the assumed bound if $d>0$, gives
\[
 Q_F\ge\min\{Q_0(n),\mathcal Q/5\},
 \quad\mathcal Q=4h^2+(6g-20n-8)h+4n^2+(6n-13)g+2g^2+14n,
\]
and both polynomials are positive in the two count ranges \chk{9.E}.
The same check gives $N_1,N_X>0$ in~\ref{compact390:middle-endpoints:M}.
\end{proof}

For $g<2n$, we use the test~\ref{comparison:linear} of Lemma~\ref{cap394:affine-overlap}.

\begin{lemma}
\label{compact390:middle-overlap}
Suppose that
\[
 B=n\ge2,\quad T=n+h,\quad n\le g\le2n-1,
 \quad h>2n,\quad h\le3g-2n+2,
\]
and let $q=n-R\ge2$.
Suppose that the converse of $H$ does not satisfy the
conditions of part~\ref{compact390:bulk-finite:1} of Theorem~\ref{compact390:bulk-finite}.
Then, with the notation above,
\[
                  C\ge0\quad\text{or}\quad P Q_F-2C B_F>0.
\]
\end{lemma}
\begin{proof}
We may assume that $C<0$. Let
\[
 V=(3g-2n-h+2)/4\ge0,\quad \Delta_0=PQ_0-2CB_0,
 \quad M=Pd-2C(g+1),\quad\Gamma=2a(3P-2C).
\]
Then $\Delta=PQ_F-2CB_F=\Delta_0+\Gamma R-M\delta$, and
$M>0$ since $d\ge4n-g>0$. Moreover $\Delta_V=\Delta_0-MV$ factorizes as
$(h-1)\mathcal D$ for a cubic $\mathcal D$ in $n,g,h$ that is positive throughout
the count range \chk{9.F}.

We claim that $\delta\le V+2R$. Otherwise, by Lemma~\ref{compact390:literal-converse}, the converse has
smaller outer crossing count $q$ and middle crossing count $g'=g-\delta<2q$,
and $S\ge4(g'+1)$, since $S-4(g'+1)=4(\delta-V)-1$ is a
nonnegative integer \chk{9.F}. At this threshold the second expression in
part~\ref{compact390:bulk-finite:1} of Theorem~\ref{compact390:bulk-finite}
is $(g'+1)(2q-g'-1)\ge0$, and it increases with $S$.
This contradicts the assumed failure of that theorem.
Finally,
\[
 \Gamma/2-M=P(4h-2n-2g-5)-2C(h-g-2)>0,
\]
so $\Delta\ge\Delta_V+(\Gamma-2M)R>0$.
\end{proof}

For $g\ge2n$, we use the test~\ref{comparison:quadratic} of Lemma~\ref{cap394:affine-overlap} when $B_F<0$.

\begin{lemma}
\label{compact390:high-caps}\label{compact390:large-middle}
Let $n\ge2$, $2\le q=n-R\le n$,
$g\ge2n$ and $g<h\le4g-2n$ be integers, and suppose that
$\delta\le U=(4g-2n-h)/5$.
Then $\max\{k-\beta_c,F_R-(1-k)\}>e+t\min(1,e/s)$ whenever
$F_Y<k$ and $\max(J_c,u,v)<k<1/2$.
\end{lemma}
\begin{proof}
Lemma~\ref{compact390:high-endpoints} gives $H_F,Q_F>0$.
If $B_F\ge0$, the affine comparison in
Lemma~\ref{cap394:affine-overlap} applies. Suppose therefore that $B_F<0$.
Then $B_U=B_0-(g+1)U\le B_F$ is negative, and
$25T_c(3/5)+30B_U>0$ \chk{9.G}, so $T_c(3/5)>0$.
For $0\le\xi\le1$, the polynomial $F$ increases with $R$ and decreases
with $\delta$, and its value $F_U$ at $R=0$, $\delta=U$ satisfies
$F_U(3/5),F_U(1)>0$ \chk{9.G}. Hence $F(3/5),F(1)>0$, and the quadratic
comparison in Lemma~\ref{cap394:affine-overlap} with $\theta=3/5$ gives
the required inequality for every $\xi>0$.
Lemma~\ref{compact390:cap-propagation} extends it throughout the region.
\end{proof}

\begin{proof}[Proof of Theorem~\ref{compact390:two-block}]
If Theorem~\ref{compact390:gap-finite} or Theorem~\ref{compact390:bulk-finite} applies to $H$ or to its converse, we are done, using Lemma~\ref{compact390:literal-converse} for the converse. Otherwise, by Lemma~\ref{compact390:count-reduction}, we may replace $H$ by its converse if necessary, so that $B=n$ and $T=n+h$ with $h>g$, and we use the bounds of that lemma. By Proposition~\ref{compact390:source-criterion}, it remains to verify its two inequalities in the crossing counts and its inequality in $u$ and $v$.

We first note that $g\ge n$. Otherwise, $g\le n-1$ and $h>g$ give $S\ge4(g+1)$, and,
at $S=4(g+1)$, the quantity $R_b$ in part~\ref{compact390:bulk-finite:1} of Theorem~\ref{compact390:bulk-finite} is $(g+1)(2n-g-1)>0$; as it increases with $S$, that theorem would apply.

Suppose first that $h\le2n$.
Part~\ref{compact390:nonnegative-slope:2} of Lemma~\ref{compact390:ratio-counts} completes this case.

Suppose next that $g<2n$ and $h>2n$. Lemma~\ref{compact390:middle-endpoints}
gives the inequalities in the crossing counts and $H_F\ge0$, $Q_F>0$,
while
Lemma~\ref{compact390:middle-overlap} gives $C\ge0$ or
$PQ_F-2CB_F>0$.
Lemma~\ref{cap394:affine-overlap}\ref{comparison:linear}, together with Lemma~\ref{compact392:cap-formulas}, then gives the inequality with right-hand side $e+t$ at $u_0$, and
Lemma~\ref{compact390:cap-propagation} extends it to the required region.

It remains to treat the case $g\ge2n$. Here the inequalities in the crossing counts follow
from $g-2n\ge0$, and the bounds of Lemma~\ref{compact390:count-reduction} give the hypotheses
of Lemma~\ref{compact390:high-caps}, which gives the inequality in $u$ and $v$.

Thus Proposition~\ref{compact390:source-criterion} applies in all cases. If we replaced $H$ by its converse, then Lemma~\ref{compact390:literal-converse} transfers the tiling back to $H$.
\end{proof}


